\chapter{Introdução}

A educação é uma preocupação central de qualquer país, e a razão é
econômica antes mesmo de ser social. A literatura de capital humano
documenta, de forma consistente ao longo de décadas, que cada ano adicional
de escolaridade se associa a um retorno privado médio de 9\% ao ano em
rendimentos futuros no mundo e de 11\% na América Latina. O maior retorno
social concentra-se justamente na etapa primária da escolarização
\cite{psacharopoulos2018returns}. Ao desenvolvimento econômico das nações,
contudo, associam-se sobretudo as habilidades cognitivas efetivamente
adquiridas pelos estudantes, medidas por avaliações padronizadas, e menos o
tempo de permanência na escola \cite{hanushek2008role}. Escolarizar sem
ensinar, adverte o Banco Mundial, é uma promessa não cumprida. Há uma
distância persistente entre \emph{schooling} e \emph{learning}, e é nela que
se decide se o investimento educacional de um país se converte em
bem-estar para as pessoas \cite{worldbank2018wdr}.

Essa preocupação é mais aguda nos países em desenvolvimento. Estimativas
conjuntas do Banco Mundial e da Organização das Nações Unidas para a
Educação, a Ciência e a Cultura (UNESCO) indicam que, após a pandemia de
Covid-19, cerca de 70\% das crianças de dez anos dos países de renda baixa
e média estavam em situação de \emph{pobreza de aprendizagem}. O termo
designa a incapacidade de ler e compreender um texto simples. Na América
Latina, a proporção estimada chegou a 80\%, ante cerca de metade antes da
pandemia \cite{worldbank2022learningpoverty}. No Brasil, a pobreza de
aprendizagem é estimada em 47\%, dos quais 45,6 pontos percentuais vêm do
déficit de \emph{aprendizagem} e apenas 2,4 do déficit de \emph{acesso} à
escola \cite{worldbank2024brazilbrief}. No ensino fundamental, o desafio
brasileiro deslocou-se, portanto, do acesso para a aprendizagem. Em países
em desenvolvimento, contexto mais próximo da realidade de parte da rede
pública brasileira, a revisão de \citeonline{glewwe2011school} encontra
efeito estatisticamente significativo para poucas características de
escolas e professores. Entre elas estão a disponibilidade de carteiras, o
conhecimento do professor sobre a matéria que ensina e a ausência docente.
Saber \emph{quais} características da escola acompanham melhores resultados
torna-se, assim, uma pergunta de alocação de recursos, além de uma
pergunta acadêmica.

A Constituição Federal, após a Emenda Constitucional
n\textordmasculine{}~59/2009, torna a educação básica obrigatória e
gratuita dos 4 aos 17 anos de idade. Ela se organiza em três etapas: a
\emph{educação infantil} (creche, para crianças de 0 a 3 anos, e
pré-escola, dos 4 aos 5), o \emph{ensino fundamental}, com duração de nove
anos e ingresso aos 6 anos de idade (Lei n\textordmasculine{}~11.274/2006),
e o \emph{ensino médio}, dos 15 aos 17. Das 47,1 milhões de matrículas da
educação básica registradas pelo Censo Escolar de 2024, 79,8\% estão na
rede pública. Nos anos iniciais do ensino fundamental, onde se situa o
2º ano, a rede pública responde por 80,4\% das matrículas, a maior parte
delas em redes municipais \cite{inep2025censo}. Compreender o que se
associa ao aprendizado na escola pública é, portanto, compreender o caso
majoritário da educação brasileira. As trajetórias que desembocam no
fundamental, porém, são desiguais desde antes dele. A pré-escola
aproxima-se da universalização (93,4\% das crianças de 4 e 5 anos), mas a
creche alcança apenas cerca de 40\% das crianças de 0 a 3 anos, com acesso
fortemente estratificado por renda. As crianças chegam, assim, ao processo
de alfabetização com pontos de partida muito distintos
\cite{ibge2025pnadeducacao}.

A alfabetização até o final do 2º ano do Ensino Fundamental, quando a
criança tem por volta de 7 a 8 anos, tornou-se uma das metas mais
consensuais da política educacional brasileira. Ela foi formalizada no
Plano Nacional de Educação e reafirmada por programas federais como o
Compromisso Nacional Criança Alfabetizada, que estabelece a meta de 80\% de
crianças alfabetizadas ao final do 2º ano até 2030. A distância até essa
meta é considerável. Na avaliação amostral do Sistema de Avaliação da
Educação Básica (Saeb) de 2023, 49,3\% dos alunos do 2º ano da rede
pública alcançaram o ponto de corte oficial de alfabetização, segundo
cálculo com os microdados e os pesos amostrais dos alunos
\cite{inep2026microdados}. A aprendizagem nos primeiros anos também se
associa à trajetória escolar. Em revisão da literatura,
\citeonline[p.~621]{hanushek2008role} relatam evidência de que, no Brasil,
habilidades cognitivas mais altas no ensino primário se associam a menos
repetência. No Egito, escolas que agregam menos aprendizagem têm mais
evasão. Investimentos nos primeiros anos de vida também podem ter retorno
elevado. Para dois programas de educação infantil da Carolina do Norte,
nos Estados Unidos, \citeonline{garcia2020quantifying} estimam taxa interna
de retorno de 13,7\% ao ano e razão benefício-custo de 7,3. Os programas
atendiam crianças de famílias desfavorecidas desde as oito semanas até os
cinco anos de idade e foram avaliados por experimento aleatorizado. O
resultado refere-se à educação infantil e serve aqui como indicação geral
do retorno de intervenções precoces. Compreender quais fatores se associam
à alfabetização nessa fase inicial é, portanto, uma pergunta tanto
acadêmica quanto diretamente relevante para a formulação de políticas
públicas.

O Saeb, conduzido pelo Instituto Nacional de Estudos e Pesquisas
Educacionais Anísio Teixeira (Inep), é o instrumento nacional que permite
formular essa pergunta com dados. Desde a edição de 2019, ele avalia
sistematicamente o 2º ano do Ensino Fundamental com uma medida de
proficiência em Língua Portuguesa (LP) estimada por Teoria de Resposta ao
Item (TRI) \cite{klein2003tri}. A partir da edição de 2023, passou a
produzir também um indicador binário de alfabetização, definido pelo Inep
por um corte de 743 pontos nessa proficiência \cite{inep2023alfabetiza}.
Paralelamente à prova aplicada ao aluno, o Saeb coleta questionários
contextuais extensos respondidos pelo professor da turma e pelo diretor da
escola. Eles cobrem temas como formação docente, práticas pedagógicas,
infraestrutura, gestão escolar e clima de segurança, e o Inep divulga
ainda o Indicador de Nível Socioeconômico (INSE) das escolas. Essa
combinação de desfechos medidos em larga escala com centenas de
covariáveis de contexto escolar permite explorar com aprendizado de máquina
supervisionado a associação entre o contexto escolar e o desempenho. O
tamanho da amostra de escolas impõe, porém, as restrições discutidas na
Seção~\ref{sec:unidade-escola}.

Há, contudo, uma lacuna específica a explorar. A literatura brasileira
sobre fatores associados ao desempenho escolar apoia-se predominantemente
em modelos de regressão multinível e de regressão quantílica
\cite{franco2007qualidade, alves2007efeito, almeida2014determinantes}. A
aplicação mais próxima de aprendizado de máquina interpretável aos
microdados do Saeb localizada nesta revisão trata do 9º ano e do Ensino
Médio \cite{tertulino2025multilevel}. No 2º ano do Ensino Fundamental, esse
instrumental permanece pouco aplicado. Trata-se de uma etapa incorporada ao
Saeb apenas a partir de 2019 e que, sem questionário respondido pelo
próprio aluno, impõe um desenho analítico distinto do adotado nas etapas
mais estudadas. Este trabalho trata dessa lacuna.

\section{Problema de pesquisa}
\label{sec:problema-pesquisa}

Este trabalho parte da seguinte pergunta central: \emph{quanto os
questionários contextuais respondidos pelo professor e pelo diretor
acrescentam, fora da amostra, à previsão da proporção de alunos
alfabetizados e da proficiência média em Língua Portuguesa das escolas
públicas com 2º ano do Ensino Fundamental, além do que o nível
socioeconômico da escola e a geografia já preveem, segundo os microdados
do Saeb 2023?}

A formulação da pergunta reflete uma característica estrutural da base. O
Saeb 2023 não aplica questionário contextual ao aluno do 2º ano. Além dos
próprios desfechos e dos identificadores de vínculo, faltam covariáveis
individuais (Seção~\ref{sec:unidade-escola}). Toda a informação preditiva
disponível é de nível turma, escola ou gestão escolar. Por isso, a escola é
adotada como unidade de análise. Os desfechos individuais são agregados por
escola, e a pergunta é enunciada no nível em que tanto os preditores
quanto as decisões de gestão educacional efetivamente operam.

\subsection{Hipóteses de pesquisa}
\label{sec:hipoteses}

Essa pergunta se desdobra em quatro hipóteses, uma para cada bloco de
variáveis: o questionário do professor, o questionário do diretor, o INSE e
a geografia. A Tabela~\ref{tab:hipoteses} as resume. Todas são enunciadas
em termos preditivos e no nível da escola. O desenho deste trabalho permite
verificar quanto cada bloco contribui para prever os desfechos agregados,
sem estabelecer relações causais (Seção~\ref{sec:registros-inferenciais}).
Cada hipótese é avaliada nas duas trilhas, pela mesma regra de decisão.
Retira-se o bloco do conjunto completo de preditores e compara-se o erro
nas mesmas dobras de uma validação cruzada repetida
(Seção~\ref{sec:robustez}). A hipótese é considerada confirmada em uma
trilha quando a retirada aumenta o MAE com intervalo de confiança de 95\%
acima de zero, pelo teste $t$ corrigido de \citeonline{nadeau2003inference},
e o aumento é de pelo menos 2\% do MAE do conjunto completo. A regra
reproduz a que foi definida antes da estimação para a seleção do modelo
final (Seção~\ref{sec:testes-significancia}). As hipóteses foram formuladas na
etapa de projeto deste trabalho, e a regra foi explicitada depois da
execução final. Reutilizar o critério já fixado para a seleção do modelo
limita a margem de escolha que essa ordem abriria.

\begin{table}[H]
    \centering
    \footnotesize
    \caption{Sub-perguntas de pesquisa e hipóteses testáveis}
    \begin{tabular}{P{0.5cm}P{3.6cm}P{5.3cm}P{2.6cm}}
        \toprule
        \textbf{\#} & \textbf{Sub-pergunta} & \textbf{Hipótese} & \textbf{Grupo temático} \\
        \midrule
        H1 & O questionário do professor acrescenta capacidade preditiva além
        do INSE, da geografia e do questionário do diretor? & Retirar o bloco
        do professor aumenta o erro de previsão. & Professor: formação,
        experiência e prática \\ \addlinespace
        H2 & O questionário do diretor acrescenta capacidade preditiva além do
        INSE, da geografia e do questionário do professor? & Retirar o bloco
        do diretor aumenta o erro de previsão. & Escola: clima, segurança e
        gestão \\ \addlinespace
        H3 & O INSE mantém capacidade preditiva própria na presença dos dois
        questionários? & Retirar o INSE aumenta o erro de previsão. & Nível
        socioeconômico \\ \addlinespace
        H4 & O contexto geográfico carrega informação própria? & Retirar as
        variáveis de região, unidade da federação e localização aumenta o
        erro de previsão. & Contexto geográfico \\
        \bottomrule
    \end{tabular}
    \vspace{0.1cm}
    \begin{flushleft}
    {\footnotesize Fonte: elaborada pelo autor.}
    \end{flushleft}
    \label{tab:hipoteses}
\end{table}

H1 reúne, no bloco do professor, formação continuada em
alfabetização e letramento, experiência docente e práticas pedagógicas em
sala de aula. H2 reúne, no bloco do diretor, clima e segurança escolar e
gestão e liderança pedagógica. Ficou fora do conjunto de hipóteses a
efetividade de ações de combate à repetência e ao abandono escolar, porque
os itens correspondentes registram apenas a \emph{existência declarada}
dessas ações, sem medir seu resultado.

H3 tem papel central. A sociologia da educação brasileira estabelece, com
razoável consenso, que o nível socioeconômico da escola é um dos preditores
mais fortes do desempenho escolar
(Seção~\ref{sec:evidencias-brasileiras}). H1 e H2 verificam se os
questionários acrescentam algo a ele, e H3 verifica se o INSE conserva
informação própria quando os questionários estão presentes. H4 aplica a
mesma lógica à geografia, ancorada nos achados de efeito-município da
literatura brasileira (Seção~\ref{sec:evidencias-brasileiras}). A
importância de cada bloco pelo SHAP (Seção~\ref{sec:xai}) é reportada ao
lado do teste, como descrição do modelo, fora da regra de decisão.

Uma limitação qualifica diretamente a Hipótese~1. O questionário do
professor do 2º ano não pergunta sobre \emph{método} de alfabetização
(fônico, silábico, global), apenas sobre práticas pedagógicas genéricas
(Seção~\ref{sec:alfabetizacao-2ef}). H1 capta, portanto, proxies de
qualidade pedagógica geral. A escolha do método de alfabetização fica fora
do alcance dos dados, e o ponto é retomado na discussão dos
resultados.

\section{Objetivos}
\label{sec:objetivos}

\subsection{Objetivo geral}

Quantificar quanto os questionários do professor e do diretor acrescentam,
fora da amostra, à previsão da proporção de alunos alfabetizados e da
proficiência média em Língua Portuguesa, além do INSE e da geografia, em
604 escolas públicas com 2º ano do Ensino Fundamental avaliadas no Saeb
2023, e hierarquizar os grupos de variáveis de maior capacidade preditiva
com modelos interpretáveis de aprendizado de máquina.

\subsection{Objetivos específicos}
\label{sec:objetivos-especificos}

\begin{enumerate}
    \item Construir uma base analítica única, com a escola como unidade de
    análise, a partir da integração de quatro fontes de microdados do Saeb
    2023 (aluno, escola, professor e diretor), com critérios de agregação e
    de seleção de variáveis documentados e reprodutíveis, livres de
    vazamento de informação e de dados pessoais sensíveis;
    \item Estimar e comparar, em duas trilhas de regressão complementares
    (proporção de alunos alfabetizados e proficiência média em Língua
    Portuguesa, ambas no nível da escola), modelos lineares regularizados
    de referência e \textit{ensembles} de árvores de decisão, avaliados por
    um protocolo experimental único com partição e validação cruzada no
    nível da escola (Seção~\ref{sec:protocolo}) e contrastados com
    \textit{baselines} simples, como a média e a mediana dos desfechos;
    \item Aplicar SHAP como método prioritário de explicabilidade ao modelo
    selecionado de cada trilha e hierarquizar a importância dos grupos
    temáticos de variáveis, usando o LIME (\textit{Local Interpretable
    Model-agnostic Explanations}) como checagem local em casos
    selecionados por regra prévia;
    \item Avaliar a robustez das conclusões frente a duas decisões
    metodológicas, a restrição da análise à rede pública e a escolha da
    escola como unidade de agregação. Para a primeira, execuções
    alternativas variam um fator de cada vez, a população analisada e o
    conjunto de preditores. Para a segunda, uma execução de sensibilidade
    repete a análise com a turma como unidade (Seção~\ref{sec:robustez});
    \item Desenvolver, como produto aplicado do trabalho, o protótipo de um
    artefato de dados (um \emph{dashboard} interativo) que traduza os
    modelos e suas explicações em um instrumento de consulta para gestores
    educacionais e potenciais financiadores de política pública.
\end{enumerate}

\section{Justificativa}
\label{sec:justificativa}

Este trabalho contribui em quatro frentes. Em primeiro lugar,
do ponto de vista metodológico, ele aplica ao Saeb uma combinação
que não foi localizada, na literatura consultada no Capítulo~2, em
trabalhos anteriores sobre o 2º ano. Trata-se de uma comparação
sistemática entre modelos lineares de referência e \textit{ensembles} de
árvore em \emph{duas} trilhas de regressão no nível da escola (proporção de
alfabetizados e proficiência média), construídas a partir de um mesmo
conjunto de preditores. Isso permite verificar se as variáveis mais
relevantes para prever a proporção de alunos alfabetizados coincidem com as
mais relevantes para prever o nível médio de proficiência. A relação desse
desenho com a regressão quantílica \cite{almeida2014determinantes} e com o
precedente de aprendizado de máquina aplicado ao Saeb
\cite{tertulino2025multilevel} é discutida na
Seção~\ref{sec:quantilica-ml}.

Em segundo lugar, do ponto de vista do rigor de dados, o trabalho registra,
para cada uma das 432 variáveis candidatas construídas a partir das quatro
bases, a decisão de inclusão ou exclusão e sua justificativa. O registro
inclui uma discussão explícita sobre vazamento de informação (a variável
alvo de alfabetização é definida por um corte na própria proficiência
contínua) e sobre o tratamento dos dados de cor/raça, fora do modelo
principal e mantidos como variável de auditoria, e de deficiência de
professores e diretores (Seção~\ref{sec:etica}).

Em terceiro lugar, do ponto de vista de política pública, hierarquizar os
grupos temáticos de contexto escolar por sua contribuição preditiva à
alfabetização, e verificar se essa contribuição se sustenta na presença do
nível socioeconômico, fornece subsídio à \emph{formulação de hipóteses} a
investigar em estudos futuros com desenho de pesquisa mais exigente. Esse
uso tem um limite. A importância preditiva em um modelo não se converte
diretamente em recomendação de política pública. A ordenação de variáveis
reflete o poder de previsão sobre os dados observados, e estimar o efeito
de intervir sobre elas exige outro desenho.

Em quarto lugar, do ponto de vista da Ciência de Dados para Negócios, o
trabalho percorre integralmente o ciclo que caracteriza a área aplicada a
um problema de decisão pública. Parte da integração de quatro bases
heterogêneas, com granularidades e chaves distintas, em uma base analítica
única no nível da escola. Organiza essa integração em um \textbf{pipeline
reprodutível}, em que os critérios de inclusão, exclusão, agregação e
codificação de cada variável residem em um módulo de configuração que
funciona como fonte única de verdade, de modo que documentação e código não
divergem. Submete os modelos a uma \textbf{comparação preditiva} sob
protocolo experimental único e traduz os modelos resultantes em
explicações auditáveis por meio de técnicas de explicabilidade. O produto
final é o protótipo de um \textbf{dashboard
interativo}, um artefato de dados nos moldes do que se apresentaria a um
comitê de investimento ou a um secretário de educação. O painel
disponibiliza essas evidências em formato explorável. Ele permite filtrar
por unidade da federação e região (variáveis de contexto geográfico
incluídas entre os preditores, Hipótese~4), consultar a ordenação dos
grupos temáticos, localizar escolas cuja proporção de alfabetizados destoa
do que seu contexto socioeconômico preveria e inspecionar a decomposição
da predição em casos concretos. A lógica de negócio subjacente é a de um
instrumento de \emph{triagem e priorização}. Com recursos escassos, o
gestor que precisa decidir onde concentrar formação docente, apoio
pedagógico ou investimento em infraestrutura dispõe de uma triagem
descritiva, inspecionável por quem a utiliza: quais escolas destoam do
esperado para seu contexto e quais grupos de fatores mais pesam na
previsão. Antes de decisões de alocação, essa triagem precisa ser combinada
com avaliações de impacto.

\section{Estrutura do trabalho}

Além desta introdução, o trabalho está organizado em mais quatro capítulos.
O Capítulo~2 apresenta a fundamentação teórica, cobrindo a avaliação do
Saeb e a Teoria de Resposta ao Item, a literatura brasileira sobre fatores
associados ao desempenho escolar, a tradição de regressão quantílica e o
aprendizado de máquina aplicado à educação, as técnicas de explicabilidade
adotadas e os fundamentos da análise com desfechos agregados no nível da
escola. O Capítulo~3 detalha a metodologia: as fontes de dados, a unidade
de análise e a delimitação da amostra, os critérios de agregação, seleção e
engenharia de variáveis, o modelo empírico e o protocolo experimental da
comparação de algoritmos, a abordagem de explicabilidade e a checagem de
robustez. O Capítulo~4 apresenta e discute os resultados: a análise
descritiva da base, o desempenho comparado dos modelos, o teste das
hipóteses pela retirada de blocos, a interpretação do modelo selecionado de
cada trilha por SHAP e LIME e a checagem de robustez. O Capítulo~5 reúne as
considerações finais, com as limitações do estudo e as possibilidades de
trabalho futuro.


\chapter{Fundamentação Teórica}

A investigação dos fatores associados ao desempenho escolar articula, ao
longo do último meio século, duas tradições de pesquisa que vêm se
aproximando nas últimas décadas. A primeira, de raiz sociológica e
econométrica, herda o desenho da função de produção educacional e da
pesquisa sobre efeito-escola. Ela busca isolar, por meio de modelos
estatísticos com pressupostos distribucionais explícitos, a contribuição
relativa de fatores familiares, escolares, docentes e de gestão sobre a
proficiência dos estudantes. A segunda, mais recente e vinculada à
mineração de dados educacionais (\textit{Educational Data Mining}) e ao
aprendizado de máquina, privilegia a capacidade preditiva e a flexibilidade
funcional de algoritmos não paramétricos. Historicamente, ela o fez ao
custo de menor transparência interpretativa, lacuna apenas parcialmente
superada pelo amadurecimento recente das técnicas de explicabilidade. Este
capítulo revisa criticamente as duas tradições e as articula ao objeto
desta pesquisa: os fatores associados à alfabetização e à proficiência em
Língua Portuguesa no 2º ano do Ensino Fundamental da rede pública
brasileira, analisados no nível da escola a partir dos microdados do Saeb
2023.

\section{Avaliação em larga escala e a mensuração do desempenho no Saeb}

O Sistema de Avaliação da Educação Básica (Saeb), conduzido pelo Instituto
Nacional de Estudos e Pesquisas Educacionais Anísio Teixeira (Inep), é a
principal fonte nacional de dados sobre desempenho escolar associados a
características contextuais de alunos, professores, diretores e escolas. A
proficiência dos estudantes não corresponde a uma nota bruta de acertos.
Ela é estimada por meio da Teoria de Resposta ao Item (TRI), que modela a
probabilidade de acerto de cada item em função de um traço latente de
habilidade e permite comparar a escala entre edições e etapas de ensino
\cite{klein2003tri}. Essa escolha metodológica tem uma consequência direta
para o desenho empírico deste trabalho, detalhada na
Seção~\ref{sec:variaveis}. Como a proficiência contínua e o indicador
binário de alfabetização são subprodutos do mesmo processo de estimação
por TRI, as demais colunas derivadas da prova do aluno (erro padrão,
respostas a itens individuais) não podem ser tratadas como preditores
independentes, sob pena de vazamento de informação.

O 2º ano do Ensino Fundamental integra o Saeb desde a edição de 2019. Em
2023, o Inep definiu o padrão oficial de alfabetização como o ponto 743 da
escala de proficiência em Língua Portuguesa \cite{inep2023alfabetiza},
critério adotado na divulgação dos resultados da edição de 2023 e
vinculado ao monitoramento do Compromisso Nacional Criança Alfabetizada. Os
pesos amostrais e os estratos dessa edição estão descritos no dicionário
que acompanha os microdados, que remete ao relatório de amostragem do Saeb
2023 \cite{inep2026microdados}. A construção do Indicador de Nível
Socioeconômico (INSE) é documentada em nota técnica própria
\cite{inep2024inse}. O INSE é usado neste trabalho como
proxy do contexto socioeconômico da escola (variável
\texttt{NIVEL\_SOCIO\_ECONOMICO}). \citeonline{costa2023reflexoes}
investigam, por meio de entrevistas com professores, como a reintrodução da
avaliação em larga escala do 2º ano foi recebida na prática pedagógica. Os
autores concluem que ainda é cedo para afirmações consistentes sobre a
intensidade dessa influência. Na literatura consultada neste trabalho, os
estudos quantitativos sobre o 2º ano são poucos em comparação com os
dedicados ao 5\textordmasculine{} e ao 9\textordmasculine{} anos.

\section{Fatores associados ao desempenho escolar: fundamentos teóricos da função de produção educacional}
\label{sec:fundamentos-teoricos}

O estudo empírico dos determinantes do desempenho escolar tem como marco
fundador o relatório coordenado por \citeonline{coleman1966equality},
encomendado pelo governo dos Estados Unidos para avaliar desigualdades
entre escolas públicas americanas. O estudo analisou dados de mais de 645
mil estudantes \cite[p.~8]{coleman1966equality} e chegou a uma conclusão
que reorientou a agenda de pesquisa em sociologia da educação. Diferenças
em infraestrutura, financiamento e currículo escolar explicavam uma
parcela surpreendentemente pequena da variação do desempenho acadêmico,
quando comparadas ao contexto socioeconômico e familiar do próprio aluno e
à composição social do corpo discente. Na prova verbal, a parcela da
variância do desempenho situada entre escolas ficou entre cerca de 5\% e
35\%, conforme a série e o grupo étnico-racial. Mais de 70\% da variação
ocorria, portanto, entre alunos da mesma escola
\cite[p.~296]{coleman1966equality}. Os fatores familiares medidos
respondiam por 10\% a 25\% da variância do desempenho individual
\cite[p.~299]{coleman1966equality}. Controlada a origem familiar, uma
medida das instalações escolares respondia sozinha por no máximo 0,02\% da
variância do desempenho verbal na 12\textordfeminine{} série de alunos
brancos e negros, e a composição do corpo discente, por 2,0\% e 6,8\%,
respectivamente \cite[p.~303]{coleman1966equality}. Esse achado inaugura a
tensão que atravessa toda a literatura subsequente sobre fatores
associados ao desempenho. É dessa tensão que parte a pergunta central
deste trabalho (Seção~\ref{sec:problema-pesquisa}): quanto o que a escola,
o professor e o diretor informam acrescenta à previsão do desempenho além
da origem socioeconômica dos alunos.

\citeonline{hanushek1970production} formaliza essa discussão por meio da
função de produção educacional, um modelo analítico que trata a escola como
um sistema em que insumos (recursos financeiros, infraestrutura,
características familiares e docentes) são transformados em resultados de
aprendizagem. Uma das descobertas mais contraintuitivas do autor é que
métricas tradicionais de qualificação docente (anos de experiência,
titulação) apresentam pouca correlação consistente com ganhos de
proficiência, ainda que a eficácia do professor varie substancialmente
entre profissionais. Essa distinção entre credenciais observáveis e
eficácia efetiva é retomada, décadas depois, por
\citeonline{hanushek2008role}. Os autores reúnem evidência de que as
habilidades cognitivas da população, medidas por avaliações padronizadas
internacionais, se relacionam fortemente com os rendimentos individuais,
com a distribuição de renda e com o crescimento econômico, e de que a
simples expansão dos anos de escolaridade não garantiu melhores condições
econômicas. Já \citeonline{glewwe2011school}, em revisão de 43 estudos de
alta qualidade sobre países em desenvolvimento publicados entre 1990 e
2010, encontram efeitos estatisticamente insignificantes para a maioria das
características de escolas e professores. As poucas exceções, como a
disponibilidade de carteiras, o conhecimento do professor sobre a matéria e
a ausência docente, oferecem, segundo os autores, pouca orientação para
políticas.

Tomados em conjunto, esses três trabalhos fornecem o andaime conceitual
(insumos familiares, insumos docentes, insumos escolares) que orienta a
seleção das variáveis candidatas deste Trabalho de Conclusão de Curso (TCC) (Seção~\ref{sec:variaveis}) e,
mais adiante, a interpretação por grupos temáticos dos valores de SHAP.

\section{Evidências brasileiras: efeito-escola, efeito-município e o papel residual do contexto escolar}
\label{sec:evidencias-brasileiras}

A aplicação da função de produção educacional ao contexto brasileiro tem em
\citeonline{barros2001determinantes} uma referência seminal. Os autores
usam dados da Pesquisa Nacional por Amostra de Domicílios (PNAD) de 1996 e
da Pesquisa sobre Padrões de Vida (PPV) de 1996/97, para jovens de 11 a 25
anos das áreas urbanas do Nordeste e do Sudeste. A escolaridade dos pais,
em especial a da mãe, é o fator mais associado aos anos de estudo
alcançados. Cada ano adicional de escolaridade dos pais corresponde a
cerca de 0,3 ano de estudo a mais para os filhos, associação mais forte que
a da renda domiciliar e superior à de três anos adicionais de escolaridade
dos professores \cite[p.~34]{barros2001determinantes}. Os insumos escolares
mantêm, ainda assim, uma margem de associação relevante. A qualidade da
infraestrutura apresenta associação tão ou mais importante que a
escolaridade docente, resultado consistente com o quadro descrito por
\citeonline{glewwe2011school} para países em desenvolvimento.
\citeonline{albernaz2002qualidade} aprofundam essa margem por meio de
modelos estatísticos multinível. Os autores estimam a associação entre a
composição socioeconômica da escola e o desempenho (efeito de pares) e
destacam a formação docente como um dos poucos insumos escolares
associados a maior equidade, além do desempenho médio.

O conceito de efeito-escola é sistematizado metodologicamente por
\citeonline{alves2007efeito}, que discutem como isolar a contribuição
própria da escola da composição socioeconômica de seus alunos. A discussão
é central para este trabalho, cujas Hipóteses~1 e~2
(Seção~\ref{sec:hipoteses}) verificam se as variáveis de professor e de
diretor acrescentam capacidade preditiva além do que o nível
socioeconômico da escola já captura. Com dados do próprio Saeb (edição de
2001, 4\textordfeminine{} série), \citeonline{franco2007qualidade} testam
quais fatores intraescolares (gestão, clima, prática docente) se associam
ao desempenho em Matemática, controlado o contexto socioeconômico. É o
precedente brasileiro metodologicamente mais próximo deste TCC, ainda que
por meio de regressão multinível clássica em vez de aprendizado de
máquina. Os autores encontram 39\% da variância da proficiência entre
escolas, ante cerca de 20\% usuais na Europa e nos Estados Unidos
\cite[p.~288]{franco2007qualidade}. Mesmo com controle pelo nível
socioeconômico médio da escola, a correção do dever de casa por todos os
professores associa-se a 12,7 pontos a mais na média da escola, e a
presença de biblioteca em sala de aula, a 4,4 pontos. Cada unidade dos
índices de equipamentos e de liderança do diretor indutora de colaboração
docente associa-se a 4,5 e 0,8 pontos, respectivamente. Parte dessas
práticas, porém, também amplia a desigualdade dentro da escola. O dever de
casa, por exemplo, eleva em 2,3 pontos o coeficiente do nível
socioeconômico do aluno \cite[p.~291]{franco2007qualidade}. O modelo
completo explica 62\% da variância das médias das escolas
\cite[p.~293]{franco2007qualidade}.

Por fim, \citeonline{soares2013efeitos} estendem a discussão ao nível do
município. Com modelos hierárquicos e dados da Prova Brasil de 2005 a 2011,
estimam efeitos de escolas e de municípios controlando as características
dos alunos e o contexto das escolas, e identificam municípios com efeitos
muito acima do esperado. O resultado sugere que o contexto geográfico e
administrativo pode carregar informação própria, não redutível ao nível
socioeconômico da escola, e é desse achado que deriva a Hipótese~4 deste
trabalho (Seção~\ref{sec:hipoteses}). O alcance dessa transposição é
limitado, porque o município é um código anonimizado nos microdados do
Saeb 2023 utilizados aqui, sem correspondência divulgada com os códigos
oficiais do Instituto Brasileiro de Geografia e Estatística (IBGE). Ele
serve de agrupamento na validação cruzada, mas não permite incorporar
covariáveis municipais nem nomear municípios, como fazem
\citeonline{soares2013efeitos} (Seção~\ref{sec:fonte-dados}).

Em conjunto, a literatura brasileira converge para um resultado de duas
camadas: o contexto socioeconômico e familiar é o preditor mais forte do
desempenho, mas subsiste uma margem residual, ainda que menor, associada a
professores, diretores e à gestão escolar e municipal. É precisamente essa
margem residual que a comparação de modelos e a interpretabilidade SHAP
deste trabalho dimensionam e hierarquizam tema a tema.

\section{Da regressão quantílica ao aprendizado de máquina: métodos quantitativos aplicados a microdados educacionais brasileiros}
\label{sec:quantilica-ml}

Um segundo eixo metodológico relevante para este trabalho é o
reconhecimento de que os fatores associados ao desempenho escolar podem
atuar de forma heterogênea ao longo da distribuição de proficiência.
\citeonline{koenker2005quantile} sistematiza a Regressão Quantílica como
alternativa à Regressão por Mínimos Quadrados Ordinários (MQO). Em vez de
estimar o efeito de covariáveis sobre a média condicional da variável
dependente, a técnica avalia esse efeito em diferentes quantis da
distribuição e revela heterogeneidades que uma análise centrada na média
esconde. Aplicada a dados educacionais, \citeonline{tian2006quantile}
encontra associações assimétricas entre variáveis de contexto familiar e o
desempenho em Matemática, mais acentuadas entre alunos situados nos
quantis inferiores da distribuição.

É nessa tradição que se insere a pesquisa de
\citeonline{almeida2014determinantes}, referência direta para este TCC.
Com os microdados da Prova Brasil 2011 e a Regressão Quantílica, o autor
identifica os determinantes dos piores e dos melhores resultados
educacionais de alunos da rede pública no 5\textordmasculine{} e no
9\textordmasculine{} anos do Ensino Fundamental. As contribuições originais
do estudo são o uso de variáveis sobre a existência de programas de
reforço escolar e a ausência de hábito de leitura (do próprio aluno e de
seus colegas de escola) e a evidência de que a associação desses fatores
com o desempenho difere entre alunos de baixo e de alto desempenho.
Programas de reforço escolar não apresentam associação consistente com o
desempenho, ao passo que a falta de hábito de leitura se associa
consistentemente a desempenho menor. O achado ilustra o tipo de
heterogeneidade que uma análise de média (MQO) não captaria.

Este trabalho compartilha com \citeonline{almeida2014determinantes} o
compromisso de ir além da média condicional, de investigar quais fatores
se associam ao desempenho e de verificar se o poder preditivo difere entre
escolas de desempenho baixo e alto. Persegue esse objetivo, porém, por uma
via metodológica distinta e complementar. Em vez de estender a regressão
quantílica a novos quantis ou a uma nova etapa de ensino, este TCC compara
modelos lineares regularizados de referência e três \textit{ensembles} de
árvore de decisão: Random Forest \cite{breiman2001random}, XGBoost e
LightGBM, estes dois implementações otimizadas do \textit{gradient
boosting} \cite{friedman2001greedy}. Os \textit{ensembles} capturam não
linearidades e interações entre as variáveis candidatas sem exigir a
especificação prévia de uma forma funcional. \citeonline{breiman2001two},
em seu influente ensaio sobre as ``duas culturas'' da modelagem
estatística, argumenta que a modelagem algorítmica, da qual os
\textit{ensembles} de árvore são exemplo, e a modelagem estatística
clássica respondem a perguntas diferentes e complementares. Essa distinção
justifica tratar os modelos lineares regularizados como \emph{modelos
lineares de referência} da comparação (Seção~\ref{sec:estrategia-empirica}),
cujos coeficientes descrevem associações condicionais
(Seção~\ref{sec:registros-inferenciais}).

Cabe distinguir, porém, o que as técnicas de explicabilidade pós-hoc
(Seção~\ref{sec:xai-revisao}) recuperam e o que não recuperam da leitura
quantílica. Regressão quantílica e SHAP respondem a perguntas distintas: a
primeira estima como a associação entre covariáveis e desfecho varia ao
longo de quantis da distribuição condicional. O segundo decompõe a predição
de um modelo já estimado em contribuições aditivas por variável, sem
qualquer pretensão de estimar parâmetros populacionais. 

Para manter a leitura por quantis, os modelos são avaliados também pela
perda de quantil (\textit{pinball loss}) em cinco pontos da distribuição
condicional dos desfechos (Seção~\ref{sec:testes-significancia}). Essa é a
função que a regressão quantílica minimiza \cite{koenker2005quantile}, e a
avaliação usa todas as escolas, sem recortar a amostra nos extremos. A
comparação permanece no plano preditivo: mede-se quão bem cada modelo prevê
o quantil $\tau$ do desfecho de cada escola, sem estimar a associação de
cada variável em cada quantil, como fazem \citeonline{almeida2014determinantes}
e \citeonline{tian2006quantile}.

Este trabalho também se insere na literatura internacional de
\textit{Educational Data Mining}, revisada por \citeonline{romero2010edm},
que situa problemas como o deste TCC (modelagem supervisionada sobre
microdados educacionais) em uma taxonomia mais ampla de tarefas de
mineração de dados aplicadas à educação. No plano internacional,
\citeonline{masci2018student} são um precedente revisado por pares
diretamente pertinente ao desenho aqui adotado. Aplicando árvores de
regressão e \textit{boosting} a dados do Programa Internacional de
Avaliação de Estudantes (PISA) de nove países, os autores modelam o
desempenho tanto de alunos quanto de \emph{escolas}. Eles mostram que os
fatores associados ao desempenho escolar agregado diferem entre sistemas
educacionais, o que indica que a análise no nível da escola é uma via
legítima e informativa para esse tipo de pergunta.

O precedente brasileiro mais direto e recente é o de
\citeonline{tertulino2025multilevel}, que integra as mesmas quatro fontes
de dados utilizadas aqui (aluno, professor, escola e diretor) em uma
comparação de algoritmos de \textit{ensemble} para classificar a
proficiência de alunos do 9\textordmasculine{} ano e do Ensino Médio. O
modelo vencedor, um Random Forest, alcança 90,2\% de acurácia e área sob a
curva ROC (AUC, do inglês \textit{Area Under the Curve}, em que ROC
abrevia \textit{Receiver Operating Characteristic}, característica de
operação do receptor) de 96,7\%. Via SHAP, os autores concluem que o nível
socioeconômico médio da escola é o preditor dominante. Para eles, ``o
desempenho acadêmico é um fenômeno sistêmico profundamente ligado ao
ecossistema da escola'' \cite[tradução nossa]{tertulino2025multilevel}.\footnote{A fonte consultada não é paginada.}
Por se apoiar em SHAP, essa leitura permanece associativa. Esse precedente
não cobre, contudo, o 2º ano do Ensino Fundamental. Nessa etapa, ao
contrário do 9\textordmasculine{} ano e do Ensino Médio, não existe
questionário do aluno, uma ausência estrutural que impõe a este trabalho
um desenho analítico distinto (Seção~\ref{sec:unidade-escola}). Essa
ausência, a unidade de análise agregada, a dupla trilha de modelagem sobre
o mesmo conjunto de preditores e a agregação de SHAP por grupo temático
formam o diferencial deste TCC frente ao precedente mais
próximo da literatura.

\section{Explicabilidade em aprendizado de máquina: abrindo a caixa-preta}
\label{sec:xai-revisao}

A adoção de \textit{ensembles} de árvore de decisão como os empregados
neste trabalho reabre, na prática, a discussão iniciada por
\citeonline{coleman1966equality} e retomada por toda a literatura sobre
fatores associados ao desempenho. Prever o desempenho escolar com acurácia
é insuficiente, e é preciso saber por quê. \citeonline{lundberg2017shap}
respondem a essa exigência unificando métodos anteriores de importância de
variáveis sob a teoria dos valores de Shapley, oriunda da teoria dos jogos
cooperativos. Disso resulta o SHAP (\textit{SHapley Additive
exPlanations}), método que atribui a cada variável uma contribuição aditiva
e consistente para a predição de cada observação. Ele permite tanto
explicações locais (por escola, no desenho deste trabalho) quanto
agregações globais (por variável ou por grupo temático de variáveis).
\citeonline{ribeiro2016lime} propõem, de forma independente, o LIME
(\textit{Local Interpretable Model-agnostic Explanations}), que aproxima o
comportamento do modelo por um modelo linear interpretável na vizinhança de
uma predição específica. Neste trabalho, o LIME serve como checagem
cruzada das explicações locais do SHAP em casos selecionados
(Seção~\ref{sec:xai}).

A escolha de empregar simultaneamente dois métodos de explicabilidade, em
vez de aceitar um único como verdade, responde à falta de consenso sobre o
que é interpretabilidade e sobre como avaliá-la, apontada por
\citeonline{doshivelez2017rigorous}. Os autores distinguem avaliações
feitas na aplicação real, com pessoas em tarefas simplificadas e por
critérios formais, sem participação humana. Este trabalho não realiza
avaliação com usuários, e a comparação entre SHAP e LIME
(Seção~\ref{sec:xai}) é feita sem participação humana, o que a situa no
terceiro tipo. \citeonline{molnar2025interpretable} consolida
esses métodos em uma síntese didática que serve de apoio complementar aos
artigos originais.

Essa literatura não é, contudo, isenta de crítica.
\citeonline{rudin2019stop} argumenta que explicações pós-hoc de modelos
caixa-preta podem ser infiéis ao comportamento real do modelo e que, em
decisões de alto risco, modelos intrinsecamente interpretáveis deveriam
ser preferidos. O desenho deste trabalho responde a essa objeção por duas
vias. Mantém modelos lineares regularizados, intrinsecamente
interpretáveis, como referência da comparação, de modo que, se o ganho
preditivo dos \textit{ensembles} não atender à regra de seleção
pré-registrada, o modelo interpretável é o selecionado
(Seção~\ref{sec:testes-significancia}). E restringe o uso do SHAP à
explicação do \emph{modelo}, sem convertê-lo em afirmação sobre o
fenômeno.

\section{Previsão, associação, explicação do modelo e inferência causal}
\label{sec:registros-inferenciais}

A articulação entre as duas tradições revisadas acima exige distinguir com
precisão quatro registros que a literatura aplicada frequentemente sobrepõe
e que este trabalho mantém separados de forma deliberada.

\textbf{Previsão} diz respeito à capacidade de antecipar o valor do
desfecho para observações não utilizadas na estimação. É o critério pelo
qual os modelos são comparados entre si (Seção~\ref{sec:protocolo}) e é
avaliada exclusivamente por desempenho fora da amostra de treino. Um modelo
pode prever bem apoiando-se em variáveis sem qualquer papel causal no
fenômeno, bastando que sejam informativas sobre ele.

\textbf{Associação} refere-se à relação estatística entre uma variável e o
desfecho, condicional às demais variáveis incluídas. É o registro em que se
enunciam as hipóteses deste trabalho (Seção~\ref{sec:hipoteses}) e o teto do
que os dados disponíveis permitem afirmar. Associações são sensíveis ao
conjunto de covariáveis incluído e não se sustentam necessariamente sob
intervenção.

\textbf{Explicação do modelo} é o que SHAP e LIME fornecem: uma decomposição
de como um modelo \emph{já estimado} chegou a determinada predição. O objeto
explicado é o modelo, não o fenômeno educacional. Duas consequências
decorrem daí: as importâncias de SHAP dependem do modelo escolhido e podem
variar entre algoritmos igualmente acurados, e uma variável pode receber
alta importância por atuar como \textit{proxy} de outra, ausente ou mal
medida na base.

\textbf{Inferência causal} exige um desenho de identificação
(aleatorização, variável instrumental, descontinuidade, diferenças em
diferenças) que este trabalho não possui, por operar sobre um corte
transversal observacional de uma única edição do Saeb. Nenhuma afirmação
causal é, portanto, formulada, e a ordenação de variáveis produzida pela
análise deve ser lida como insumo para a formulação de hipóteses causais a
serem testadas em estudos com desenho apropriado, não como estimativa do
efeito de intervir sobre essas variáveis.

Métodos recentes de aprendizado de máquina para efeitos causais não
dispensam essa exigência. As florestas causais de
\citeonline{wager2018causal} estimam efeitos heterogêneos de um tratamento,
e o \textit{double machine learning} de \citeonline{chernozhukov2018double}
usa algoritmos como \textit{random forests}, lasso e \textit{boosting} para
controlar covariáveis na estimação de efeitos médios. Os dois flexibilizam
a forma funcional, mas dependem de um tratamento definido e de uma
hipótese de identificação. Em geral, é a de não confundimento (condicional
às covariáveis, a atribuição do tratamento independe dos resultados
potenciais), acompanhada de sobreposição entre tratados e não tratados, ou
a de uma variável instrumental. A biblioteca CausalML, que reúne métodos de
\textit{uplift} e de efeitos heterogêneos, recomenda aplicá-los a dados
experimentais e adverte que, em dados observacionais, confundidores
omitidos enviesam a estimativa \cite{chen2020causalml}. O desenho deste
trabalho não satisfaz essas condições. Seus preditores são centenas de
itens de questionário correlacionados, sem um tratamento definido,
observados em um único corte transversal. Além disso, a composição social
das escolas torna implausível supor que as práticas declaradas por
professores e diretores independam de fatores não observados. Esses
métodos são um passo para estudos futuros que isolem uma prática
modificável apontada por esta análise e disponham de variação capaz de
identificá-la, como várias edições do Saeb ou um desenho quase
experimental.

\textbf{Regra de redação.} Os quatro registros resultam em uma regra
única para todo o trabalho. Os resultados são enunciados como associações
preditivas no nível da escola (``escolas com maior X apresentam, em média,
maior proporção de alfabetizados''), nunca como efeitos causais nem como
afirmações sobre alunos individuais. Os valores de SHAP e LIME descrevem o
que o modelo aprendeu com os dados de treinamento, e a ordenação de
variáveis serve como ponto de partida para hipóteses causais a testar em
estudos com desenho apropriado. O
fundamento da restrição ao nível da escola está na
Seção~\ref{sec:nivel-escola-teoria}.

\section{Análise no nível da escola: desfechos agregados, proporções e o alcance ecológico das conclusões}
\label{sec:nivel-escola-teoria}

A adoção da escola como unidade de análise insere este trabalho em uma
literatura metodológica própria, com dois fundamentos e uma advertência.

O primeiro fundamento é econométrico: desfechos expressos como proporções
(caso da proporção de alunos alfabetizados de cada escola) são
variáveis limitadas ao intervalo $[0,1]$, para as quais a literatura
recomenda especificações que respeitem esse suporte.
\citeonline{papke1996fractional} propõem o modelo de resposta fracionária
(\textit{fractional logit}), que modela a esperança condicional da
proporção por meio de uma função de ligação logística sem exigir
transformações ad hoc dos valores extremos 0 e 1.
\citeonline{ferrari2004beta} desenvolvem alternativa paramétrica pela
regressão beta. Neste trabalho, essas especificações informam o papel do
modelo linear de referência da trilha de proporção
(Seção~\ref{sec:estrategia-empirica}), enquanto os \textit{ensembles} de
árvore, por não imporem forma funcional, acomodam naturalmente o desfecho
fracionário, com a checagem de que as predições respeitem o intervalo
unitário.

O segundo fundamento é de desenho. Quando todos os preditores disponíveis
são medidos no nível do grupo, replicá-los para cada indivíduo do grupo
apenas repete a mesma linha muitas vezes, sem acrescentar informação. O
fenômeno, conhecido como pseudorreplicação, inflaciona artificialmente o
tamanho aparente da amostra e o desempenho aparente dos modelos. Agregar o
desfecho ao nível em que os preditores variam torna explícita a estrutura
de informação dos dados, em vez de escondê-la sob dezenas de milhares de
linhas redundantes.

A advertência é clássica: \citeonline{robinson1950ecological} demonstrou
que correlações calculadas sobre agregados (as chamadas correlações
ecológicas) podem diferir substancialmente (em magnitude e até em sinal)
das correlações no nível individual, erro de transposição conhecido
como \emph{falácia ecológica}. Por isso, a regra de redação deste
trabalho (Seção~\ref{sec:registros-inferenciais}) mantém as associações
estimadas no nível da escola.

\section{Alfabetização no 2º ano do Ensino Fundamental: especificidades do objeto}
\label{sec:alfabetizacao-2ef}

A alfabetização, tal como tratada neste trabalho, não é um constructo
simples. \citeonline{soares2017alfabetizacao} a define como um processo
complexo, composto por três facetas (linguística, interativa e
sociocultural), e reconstrói o embate histórico entre métodos de
alfabetização analíticos e sintéticos (fônico, silábico, global) que ainda
hoje organiza o debate público e de política educacional no Brasil. A
evidência experimental internacional ajuda a qualificar esse embate. A
síntese do National Reading Panel, que reuniu experimentos sobre o ensino
inicial da leitura em língua inglesa, concluiu que a instrução fônica
explícita e sistemática produz ganhos de leitura maiores que programas sem
instrução fônica ou com instrução fônica assistemática. O efeito médio é
moderado (\textit{d} = 0,44) e maior quando o ensino começa na educação
infantil ou no 1º ano (\textit{d} = 0,56 e 0,54) do que a partir do 2º ano
(\textit{d} = 0,27). A mesma síntese não encontrou diferença estatística
entre os tipos de programa fônico sistemático, inclusive os sintéticos
\cite[p.~2-92 e 2-93]{nrp2000}. Há suporte, portanto, para afirmar que o
ensino sistemático das relações entre letras e sons importa na etapa
estudada aqui. A escolha entre métodos fônicos específicos não tem a mesma
sustentação, e a transposição desses resultados, obtidos em língua
inglesa, para o português requer cautela. Essa
discussão expõe uma limitação estrutural deste TCC. O questionário do
professor do 2º ano
do Ensino Fundamental no Saeb 2023 não pergunta sobre o método de
alfabetização adotado em sala de aula, apenas sobre práticas pedagógicas
de natureza mais genérica. Consequentemente, as variáveis de prática
docente utilizadas neste trabalho (Seção~\ref{sec:variaveis}) funcionam
como \textit{proxies} de qualidade pedagógica geral, não como uma medida da
escolha de método de alfabetização especificamente, o que qualifica de
forma importante a interpretação da Hipótese~1 (Seção~\ref{sec:hipoteses}).
Some-se a isso a constatação de \citeonline{costa2023reflexoes} de que a
própria avaliação em larga escala do 2º ano ainda é
recente na prática das escolas.

\section{Síntese e posicionamento da pesquisa}

A fundamentação deste capítulo permite situar a contribuição deste TCC. Da função de produção educacional
(Seção~\ref{sec:fundamentos-teoricos}) e das evidências brasileiras de
efeito-escola e efeito-município (Seção~\ref{sec:evidencias-brasileiras}),
herda-se a expectativa teórica de que o contexto socioeconômico domina a
associação com o desempenho escolar sem esgotá-la. Resta uma margem
empírica para professores, diretores e gestão municipal, que a comparação
de modelos e a interpretabilidade SHAP deste trabalho dimensionam. Da tradição de regressão quantílica
(Seção~\ref{sec:quantilica-ml}), herda-se o compromisso de investigar
heterogeneidade além da média condicional, perseguido aqui por duas vias:
a comparação entre modelos lineares de referência e \textit{ensembles} de
árvore, interpretada por SHAP (Seção~\ref{sec:xai-revisao}), e a avaliação
desses modelos em vários quantis pela perda de quantil
(Seção~\ref{sec:testes-significancia}). Da literatura de desfechos
agregados (Seção~\ref{sec:nivel-escola-teoria}), herda-se a especificação
adequada a proporções e o limite ecológico das conclusões. Do precedente
mais próximo em aprendizado de máquina aplicado ao Saeb
\cite{tertulino2025multilevel}, herda-se o desenho de integração de quatro
fontes de dados. Identifica-se ali, porém, uma lacuna estrutural
específica, a ausência de questionário do aluno no 2º ano do Ensino
Fundamental (Seção~\ref{sec:alfabetizacao-2ef}). Esta pesquisa a endereça
com um desenho necessariamente distinto, detalhado no capítulo de
Metodologia.

\chapter{Metodologia}

Este capítulo descreve, de forma replicável, os procedimentos empíricos
que respondem à pergunta de pesquisa. As Seções~\ref{sec:fonte-dados}
a~\ref{sec:variaveis} tratam dos dados, da unidade de análise, dos pesos,
do viés associado ao INSE e das variáveis. As
Seções~\ref{sec:estrategia-empirica} a~\ref{sec:robustez} tratam da
comparação de modelos, da seleção do modelo final, da explicabilidade (XAI,
do inglês \textit{eXplainable Artificial Intelligence}), da robustez e do
valor incremental dos blocos.

\section{Fonte de dados e integração das bases}
\label{sec:fonte-dados}

\subsection{O Saeb 2023 e o 2º ano do Ensino Fundamental}

O Saeb 2023 disponibiliza, para o 2º ano do Ensino
Fundamental, seis arquivos de microdados: \texttt{TS\_ALUNO\_2EF},
\texttt{TS\_ESCOLA}, \texttt{TS\_PROFESSOR}, \texttt{TS\_DIRETOR},
\texttt{TS\_SECRETARIO\_\allowbreak MUNICIPAL} e \texttt{TS\_ITEM}.
Os dois últimos foram excluídos do escopo.
\texttt{TS\_\allowbreak SECRETARIO\_MUNICIPAL} tem granularidade municipal,
incompatível com o nível de análise deste trabalho. Além disso, o
identificador de município do Saeb é um código anonimizado, sem
correspondência divulgada com os códigos oficiais do IBGE, o que impede
vincular a base a fontes externas de dados municipais e restringe o uso do
município ao papel de variável de agrupamento na validação cruzada
(Seção~\ref{sec:protocolo}). Já \texttt{TS\_ITEM} refere-se à análise
psicométrica, via TRI, de funcionamento diferencial de item (DIF, do inglês
\textit{Differential Item Functioning}), uma pergunta de pesquisa distinta
da investigação de fatores contextuais associados ao desempenho aqui
proposta. As quatro bases restantes (aluno, escola, professor e diretor)
são integradas em uma única base analítica, descrita a seguir.

\subsection{Chaves, filtros e registros sem correspondência}

\textbf{Chaves e filtros de integração.} A junção das quatro fontes segue
chaves e filtros distintos, resumidos na Tabela~\ref{tab:fontes-dados}. O
questionário do professor é unido pela chave \texttt{ID\_TURMA},
restrito a \texttt{ID\_SERIE~==~2} e \texttt{SQ\_QUESTIONARIO~==~5}, código
que o dicionário de dados do Inep identifica como o bloco \emph{Professor
Único, 2º e 5º ano}. O primeiro filtro seleciona os registros referentes ao
2º ano. O segundo seleciona, entre os múltiplos blocos de
questionário que a base do professor concatena em um mesmo arquivo, aquele
efetivamente aplicado ao docente regente da turma avaliada (sem ele, um
mesmo professor apareceria replicado em mais de um registro por turma, com
conjuntos distintos de itens respondidos). A contagem de registros
por \texttt{ID\_TURMA} confirmou a relação um para um entre professor e
turma resultante desses filtros, com 1.045 questionários para 1.045
turmas.

O questionário do diretor é unido pela chave \texttt{ID\_ESCOLA}, também
restrito a \texttt{ID\_SERIE~==~2}, em relação um-para-um com a escola, e
cobre 95,4\% das escolas da base analítica. Já \texttt{TS\_ESCOLA}, única
fonte do Indicador de Nível Socioeconômico (INSE), cobre exclusivamente a
rede pública. Das 1.107 escolas com alunos de proficiência válida, 764 são
públicas, e 704 delas dispõem do INSE (92,1\% das públicas e 63,6\% do
total). As implicações dessa assimetria de cobertura são discutidas na
Seção~\ref{sec:vies-inse}.

\textbf{Registros sem correspondência.} A integração parte da base de
alunos, por meio de junção à esquerda. Um aluno cuja turma não possua
questionário de professor correspondente, ou cuja escola não possua
questionário de diretor, permanece na base com as variáveis daquele bloco
ausentes, sendo tratado pelo procedimento de imputação descrito na
Seção~\ref{sec:protocolo}. Escolas que percam a totalidade dos blocos de
preditores são excluídas, por não restar informação explicativa. A
Tabela~\ref{tab:cobertura} mostra quantos alunos, turmas e escolas
permanecem em cada etapa de filtragem e junção.

\begin{table}[H]
    \centering
    \footnotesize
    \caption{Fontes de microdados do Saeb 2023 utilizadas e critérios de integração}
    \begin{tabular}{P{2.9cm}P{1.6cm}P{2.3cm}P{2.9cm}P{3.2cm}}
        \toprule
        \textbf{Fonte} & \textbf{Unidade de observação} & \textbf{Chave de junção} & \textbf{Cobertura} & \textbf{Papel na base analítica} \\
        \midrule
        TS\_ALUNO\_2EF & Aluno & --- & Universo do 2º ano & Desfechos (agregados por escola) e identificadores de vínculo \\ \addlinespace
        TS\_PROFESSOR & Turma & ID\_TURMA & Universo (público+privado) & Preditores de perfil e prática docente (agregados por escola) \\ \addlinespace
        TS\_DIRETOR & Escola & ID\_ESCOLA & 95,4\% das escolas da base & Preditores de gestão escolar \\ \addlinespace
        TS\_ESCOLA & Escola & ID\_ESCOLA & Só rede pública (92,1\% das escolas públicas) & Única fonte do INSE \\
        \bottomrule
    \end{tabular}
    \vspace{0.1cm}
    \begin{flushleft}
    {\footnotesize Fonte: elaborada pelo autor, com base no dicionário de
    dados do Saeb 2023 (Inep).}
    \end{flushleft}
    \label{tab:fontes-dados}
\end{table}

\texttt{TS\_ALUNO\_2EF} não contribui com covariáveis de
contexto do aluno: no 2º ano não há questionário
respondido pelo estudante, de modo que essa base fornece apenas os
desfechos, os identificadores de vínculo com turma e escola e as variáveis
de estrato e localização.

\begin{table}[H]
    \centering
    \footnotesize
    \caption{Alunos, turmas e escolas em cada etapa de filtragem e integração}
    \begin{tabular}{L{7.2cm}C{1.9cm}C{1.9cm}C{1.9cm}}
        \toprule
        \textbf{Etapa} & \textbf{Alunos} & \textbf{Turmas} & \textbf{Escolas} \\
        \midrule
        Registros iniciais (TS\_ALUNO\_2EF) & 37.104 & 1.612 & 1.181 \\
        Consistência com o Censo Escolar & 36.989 & 1.611 & 1.181 \\
        Prova de LP preenchida & 28.532 & 1.504 & 1.107 \\
        Proficiência em LP calculável (TRI) & 28.516 & 1.504 & 1.107 \\
        Rede pública (escopo principal) & 20.939 & 1.069 & 764 \\
        Com INSE (junção com TS\_ESCOLA) & 19.361 & 985 & 704 \\
        \quad dos quais, turma com questionário de professor & 14.613 & 730 & 527 \\
        \quad dos quais, escola com questionário de diretor & 16.466 & 825 & 587 \\
        Escola com ao menos um bloco de preditores & 17.269 & 867 & 616 \\
        \textbf{Escola com 10 ou mais alunos (base analítica final)} & 17.183 & 854 & 604 \\
        \midrule
        Coorte interna (professor e diretor, sensibilidade) & 13.810 & 688 & 498 \\
        \bottomrule
    \end{tabular}
    \vspace{0.1cm}
    \begin{flushleft}
    {\footnotesize Fonte: elaborada pelo autor, com base nos microdados do Saeb
    2023 (Inep). As linhas recuadas não são filtros. Elas indicam quantos
    registros da etapa anterior têm o questionário correspondente, já que a
    integração usa junção à esquerda.}
    \end{flushleft}
    \label{tab:cobertura}
\end{table}

\subsection{Filtros de elegibilidade e delimitação do escopo}

A coorte analítica principal é delimitada por três filtros sucessivos no
nível do aluno, aplicados \emph{antes} da agregação: consistência com o
Censo Escolar, preenchimento adequado do questionário e proficiência em
Língua Portuguesa calculável pela TRI. Como \texttt{TS\_ESCOLA} é a única
fonte do INSE e cobre apenas a rede pública, o escopo principal deste
trabalho é adicionalmente restrito a \texttt{IN\_PUBLICA~==~1}. Para
avaliar se essa restrição altera as conclusões, execuções alternativas
descritas na Seção~\ref{sec:robustez} repetem a trilha da proporção de
alfabetizados sobre o universo de escolas públicas e privadas, sem a
variável de INSE.

\section{Unidade de análise: da observação-aluno à observação-escola}
\label{sec:unidade-escola}

A decisão metodológica mais consequente deste projeto é a adoção da
\textbf{escola como unidade de análise}: cada linha da base analítica
corresponde a uma escola pública com 2º ano avaliado no Saeb 2023, os
desfechos são agregados dos alunos para a escola e os preditores são
consolidados no mesmo nível. A seguir vêm as razões da decisão, o
procedimento de agregação e suas implicações, inclusive os custos.

\subsection{Por que a escola é a unidade de análise}
\label{sec:limitacao-estrutural}

No 2º ano do Ensino Fundamental não há questionário contextual respondido
pelo aluno. As variáveis preditoras candidatas originam-se inteiramente de
questionários de turma (professor) e de escola (diretor, escola,
geografia), e alunos de uma mesma turma compartilham exatamente o mesmo
vetor de variáveis preditoras. Uma \textbf{análise exploratória} da base
integrada dimensiona essa característica. Na coorte de escolas públicas
com questionários de professor e de diretor, 13.810 alunos distribuem-se
por 688 turmas em 498 escolas. Na base analítica final, com junção à
esquerda, são 17.183 alunos em 604 escolas. A correlação intraclasse (ICC,
do inglês \textit{intraclass correlation coefficient}) da proficiência é
próxima de 0,20, o que situa o tamanho amostral \emph{efetivo} na ordem de
poucos milhares de observações-equivalentes, muito abaixo do número de
linhas-aluno.

Manter o aluno como unidade de observação, nessas condições, significa repetir
a mesma linha de preditores dezenas de vezes, variando apenas o desfecho:
uma estrutura de \emph{pseudo-replicação} que inflaciona o tamanho
aparente da amostra sem acrescentar informação explicativa
(Seção~\ref{sec:nivel-escola-teoria}). O mesmo ICC impõe, ainda, um
\textbf{teto estrutural} à modelagem no nível individual: se apenas cerca
de 20\% da variância da proficiência reside \emph{entre} escolas, um
modelo alimentado apenas por preditores de escola não pode explicar mais
que essa fração da variância individual, e preditores de turma ficam
limitados, analogamente, à parcela da variância situada entre turmas. A agregação reconhece esse limite e reformula o alvo para o
nível em que a variância explicável de fato reside. Além do argumento
estatístico, há um argumento de decisão: a escola é a menor unidade sobre a
qual a gestão educacional efetivamente atua (alocação de professores,
formação, infraestrutura), de modo que desfechos e explicações no nível da
escola são diretamente acionáveis pelo usuário final do
dashboard (Seção~\ref{sec:justificativa}).

\subsection{Agregação dos desfechos e dos preditores}
\label{sec:agregacao}

\textbf{Desfechos.} Os dois desfechos individuais são agregados por escola:
o indicador binário de alfabetização torna-se a \textbf{proporção de alunos
alfabetizados} da escola,
\begin{equation}
    Y_e^{\text{alf}} = \frac{1}{n_e} \sum_{i \in e} \mathbf{1}\left[\text{PROFICIENCIA\_LP\_SAEB}_i \geq 743\right],
    \label{eq:prop-alfabetizado}
\end{equation}
em que $n_e$ é o número de alunos avaliados da escola $e$. A proficiência
contínua torna-se, por sua vez, a \textbf{proficiência média} da escola,
$Y_e^{\text{prof}} = \frac{1}{n_e} \sum_{i \in e}
\text{PROFICIENCIA\_LP\_SAEB}_i$. O número de alunos avaliados $n_e$ é
retido como variável auxiliar. Escolas com $n_e$ muito pequeno produzem
proporções instáveis, e por isso a base exige ao menos 10 alunos avaliados
por escola (Tabela~\ref{tab:cobertura}). A Seção~\ref{sec:robustez} repete
a análise com o corte de 20 alunos.

\textbf{Preditores.} As variáveis de diretor, de escola e de geografia já
são de nível escola e não requerem transformação. As variáveis do
questionário do professor são de nível turma. Quando uma escola possui
mais de uma turma de 2º ano avaliada, elas são consolidadas por escola:
variáveis numéricas e ordinais pela média, variáveis binárias pela
proporção de turmas com resposta afirmativa e variáveis nominais pela moda,
com o número de turmas da escola retido como variável de contexto. Essa
consolidação não é inócua. Na base final, 41,4\% das escolas têm mais de
uma turma avaliada, e a agregação por escola funde vetores docentes
distintos. Por isso, além da retenção do
número de turmas, a Seção~\ref{sec:robustez} traz uma execução de
sensibilidade com a \emph{turma} como unidade de agregação, que preserva o
bloco docente intacto. Os critérios de agregação de cada variável estão
registrados no mesmo módulo de configuração que documenta sua inclusão
(Seção~\ref{sec:variaveis}).

\subsection{Implicações da agregação}
\label{sec:implicacoes-agregacao}

Três implicações alteram a leitura de todos os resultados.

\textbf{Queda do tamanho amostral e das métricas.} A base analítica passa
de dezenas de milhares de linhas-aluno para algumas centenas de
linhas-escola (604 na base final). Isso impõe uma consequência direta à
modelagem: com 363 preditoras candidatas para algumas centenas de
observações, a regularização dos modelos e a verificação explícita do valor
incremental de cada bloco de variáveis (Seção~\ref{sec:robustez}) tornam-se
etapas obrigatórias, sob pena de sobreajuste. Além disso,
espera-se que as métricas de desempenho reportadas sejam
\emph{substancialmente inferiores} às que um protocolo com partição
aleatória por aluno produziria, e isso não é uma deficiência do desenho,
mas sua correção: parte relevante do desempenho aparente sob partição por
aluno decorre de o modelo reencontrar no teste vetores de preditores que já
viu no treino. As métricas obtidas no nível da escola aproximam a
capacidade de generalização para outras escolas da mesma população e da
mesma edição, que é a pergunta relevante para o uso do
instrumento. Elas não garantem desempenho semelhante em edições futuras do
Saeb.

\textbf{Perda da leitura individual.} Ao agregar, o trabalho abdica de
qualquer afirmação sobre alunos individuais. Não há predição de risco
individual de não alfabetização, e as associações estimadas valem para
escolas, não para crianças (falácia ecológica,
Seção~\ref{sec:nivel-escola-teoria}). O que se ganha em contrapartida é um
desfecho contínuo mais informativo que o binário individual (a proporção
carrega intensidade, não apenas estado) e diretamente alinhado às metas
públicas, que são formuladas exatamente nesses termos: a proporção de
crianças alfabetizadas é o indicador oficial do Compromisso Nacional
Criança Alfabetizada.

\textbf{Simplificação do protocolo de validação.} Com uma linha por escola,
a partição treino/teste e a validação cruzada operam diretamente sobre
escolas, o que elimina por construção o vazamento entre colegas de turma.
Permanece
uma dependência residual possível entre escolas de um mesmo município ou
rede (gestão municipal comum, Seção~\ref{sec:evidencias-brasileiras}). A
análise de sensibilidade examina a estabilidade das conclusões sob
validação cruzada agrupada por município.

\section{Pesos amostrais e alcance das conclusões}
\label{sec:pesos}

O Saeb 2023 é uma avaliação amostral no 2º ano do Ensino Fundamental. Cada
escola sorteada representa outras escolas da rede, e o peso amostral indica
quantas, conforme o desenho descrito nos microdados \cite{inep2026microdados}.
Esse peso difere do peso de precisão $n_e$, o número de alunos avaliados da
escola, tratado no fim desta seção.

A necessidade de ponderar depende do que se quer estimar.
\citeonline{solon2015weighting} distinguem dois propósitos. Para estimar
estatísticas descritivas de uma população, como a proporção de alunos
alfabetizados da rede, a ponderação é necessária sempre que a amostra não
representa essa população. Para estimar relações entre variáveis, o peso
pode ser desnecessário à consistência quando a probabilidade de seleção
depende apenas das variáveis explicativas, e ponderar pode então reduzir a
precisão. Os autores recomendam reportar as estimativas com e sem peso,
porque o contraste entre elas serve de diagnóstico. O estimando deste
trabalho pertence ao segundo tipo: os padrões preditivos e a importância das
variáveis \emph{no conjunto de escolas analisado}. Os modelos não produzem
estimativas populacionais do 2º ano do Brasil, e nenhuma generalização
automática para o país é formulada a partir deles.

Os pesos variam muito na população principal. O coeficiente de variação é
de 2,14, e o efeito de desenho decorrente dessa variação,
$\text{deff} = n\sum_e w_e^2 \big/ \left(\sum_e w_e\right)^2$, é de 5,59.
Para fins de precisão, as 604 escolas equivalem a cerca de 108 escolas
efetivas. Os pesos também se associam, ainda que fracamente, aos desfechos
(correlação de Spearman de 0,15 com a proporção de alfabetizados e de 0,13
com a proficiência média), e por isso as médias ponderadas superam as não
ponderadas. A proporção média de alfabetizados passa de 43,9\% para 52,0\%
(18,4\% a mais), e a proficiência média, de 731,6 para 739,9 pontos (1,1\%
a mais). Estatísticas descritivas sem peso subestimariam, portanto, o
desempenho da rede, e toda média que se refira à rede é reportada com peso.
Em 4,6\% das escolas, as que não têm questionário do diretor, o peso é imputado a
partir do peso médio de seus alunos, ajustado pela razão mediana entre os
dois pesos.

A estratégia segue a recomendação de reportar as duas versões. Os modelos
são treinados sem o peso amostral, e as métricas de erro são reportadas sem
peso e ponderadas por ele. Na análise de sensibilidade, o treino com peso
piora o erro ponderado nas duas trilhas. O ganho sobre o \textit{baseline},
medido com peso, cai de 15,3\% para 4,6\% na trilha de proporção e de
15,4\% para 9,7\% na de proficiência, e truncar os pesos no percentil 99
recupera só parte da perda (9,6\% e 9,5\%). O resultado é coerente com a
perda de precisão esperada quando os pesos variam tanto: o modelo treinado
sem peso prevê melhor a própria rede ponderada. O mesmo contraste vale para
a importância das variáveis. Calculada nas 604 escolas com o modelo
selecionado (Seção~\ref{sec:xai}), a participação de cada tema no SHAP
produz praticamente o mesmo ranking sem e com peso. A correlação de
Spearman entre os dois rankings é de 0,998 na trilha de proporção e de
0,994 na de proficiência, com limite inferior do intervalo \textit{bootstrap}
de 95\% de 0,992 e 0,989. O primeiro tema é o mesmo nas duas versões, e
coincidem quatro e cinco dos cinco primeiros temas, respectivamente.

O peso de precisão $n_e$ responde a outra questão. Como a proporção de
alfabetizados de uma escola pequena oscila mais por acaso, todos os modelos
da trilha de proporção usam $n_e$ como peso no ajuste. As métricas
principais tratam cada escola com o mesmo peso, coerente com a unidade de
decisão, e o MAE ponderado por $n_e$ é reportado como complemento
(Seção~\ref{sec:estrategia-empirica}).

\section{Viés de seleção associado à disponibilidade do INSE}
\label{sec:vies-inse}

A restrição do escopo principal às escolas com INSE disponível não é
aleatória: ela seleciona a rede pública e, dentro dela, 91,8\% das escolas
com pelo menos 10 alunos avaliados (604 de 658). Se as escolas com INSE
disponível diferirem sistematicamente
das demais em características associadas ao desfecho (porte, localização,
dependência administrativa, região), as associações estimadas sobre a
subamostra não se transferem diretamente ao universo.

Para dimensionar esse risco, a análise faz uma \textbf{comparação
descritiva entre as escolas incluídas e as excluídas} por falta do INSE,
entre as escolas públicas com pelo menos 10 alunos avaliados. Os dois
grupos são contrastados nos desfechos agregados, no porte, no número de
turmas, na localização (rural e capital), na região e na presença dos
questionários de professor e de diretor, pela diferença média padronizada e
pelo teste de Mann-Whitney. O resultado aparece no Capítulo~4 e qualifica o
alcance das conclusões, junto com as execuções de robustez da
Seção~\ref{sec:robustez}, que prescindem do INSE e operam sobre o universo
completo.

\section{Variáveis}
\label{sec:variaveis}

\subsection{Variáveis-alvo e o problema do vazamento de informação}

Este trabalho define duas variáveis-alvo complementares no nível da escola,
ambas construídas a partir da proficiência contínua em Língua Portuguesa
(\texttt{PROFICIENCIA\_LP\_SAEB}), estimada pela TRI \cite{klein2003tri}. A
primeira é a \textbf{proporção de alunos alfabetizados} da escola
(Equação~\ref{eq:prop-alfabetizado}), em que o indicador individual de
alfabetização é definido pelo Inep (\citeyear{inep2023alfabetiza}) como
$\mathbf{1}[\text{PROFICIENCIA\_LP\_SAEB}_i \geq 743]$. A segunda é a
\textbf{proficiência média} da escola. Como ambas derivam do mesmo processo
de estimação por TRI, todo o bloco de colunas da tabela do aluno derivado
desse processo constitui vazamento de informação e foi excluído do
conjunto de preditores nas duas trilhas. O bloco inclui a própria
proficiência, quando não compõe o alvo da trilha, o erro-padrão de
estimação e as respostas a itens individuais da prova. Pela mesma razão,
os dois alvos não podem figurar um como preditor do outro, porque são
transformações distintas da mesma variável de origem.

\subsection{Seleção, codificação e engenharia das variáveis preditoras}

As quatro bases integradas somam 591 colunas, a partir das quais se
constroem 432 variáveis candidatas. Cada variável candidata foi avaliada
individualmente quanto à sua inclusão como preditor, com a decisão e a
justificativa registradas em um módulo de configuração que funciona como
fonte única de verdade para todo o pipeline (\texttt{dicionario.py},
exportado como tabela versionada a cada execução). Isso evita divergência
entre o código de modelagem e a documentação da base. Além do vazamento
descrito acima, foram excluídas as colunas de texto aberto, as colunas com
mais de 30\% de ausência na amostra filtrada e a variável
\texttt{ESTRATO}, um metadado do desenho amostral \cite{inep2026microdados}.
Em testes preliminares, \texttt{ESTRATO} se comportava como
quase-identificador de escola e dominava desproporcionalmente a importância
agregada de SHAP, sem corresponder a um fator de contexto substantivo. O
processo resultou em 363 variáveis preditoras. Cada uma foi codificada como
binária, ordinal (tipo Likert), numérica ou nominal a partir da coluna de
tipo do dicionário do Inep e do texto das opções de resposta (nunca da
posição da letra) e classificada em um de dezenas de grupos temáticos (por
exemplo, ``práticas pedagógicas em sala de aula'' ou ``clima escolar''). Esses
grupos permitem agregar a importância de SHAP por tema, além da leitura por
coluna isolada (Seção~\ref{sec:xai}). Nas escalas que incluem opções
alheias à escala (``Não tem'', ``Não uso'', ``Não se aplica'', ``Não sei''),
essas opções geram um indicador binário próprio, em vez de serem ordenadas
artificialmente junto aos demais níveis. O mesmo módulo de configuração
registra, para cada variável de nível turma, a regra de agregação por
escola descrita na Seção~\ref{sec:agregacao}.

Duas decisões complementam esse conjunto. Primeiro, foram retirados da
especificação principal dois itens que medem a percepção do respondente
sobre o próprio desfecho: a concordância do professor e do diretor com a
afirmação de que a maioria dos estudantes apresenta problemas de
aprendizagem (\texttt{PROF\_TX\_Q080} e \texttt{DIR\_TX\_Q030}). Como a
resposta tende a refletir o desempenho já observado, esses itens carregam
risco de causalidade reversa e transformariam a ordenação de importâncias
em um termômetro do próprio alvo. Uma análise de sensibilidade os
reintroduz para dimensionar seu efeito (Seção~\ref{sec:robustez}). Segundo, as escolas sem
questionário apresentam desfechos médios um pouco menores que as demais,
embora as diferenças não alcancem significância estatística ao nível de
5\%. Como precaução, foram incluídos dois
indicadores binários de escola sem questionário de professor e sem
questionário de diretor, de modo que o modelo distinga a não resposta de
uma resposta mediana, que é o valor atribuído pela imputação.

\subsection{Considerações éticas e delimitação de escopo}
\label{sec:etica}

\textbf{Natureza dos dados utilizados.} Os microdados do Saeb 2023 são
\emph{públicos}, disponibilizados pelo Inep em portal de dados abertos, e
distribuídos já \emph{anonimizados}: não contêm nome, documento ou qualquer
identificador direto de aluno, professor ou diretor, e os identificadores de
turma e escola são códigos sem correspondência divulgada com pessoas
naturais. Seu uso neste trabalho tem \emph{finalidade exclusivamente
acadêmica}, em conformidade com os termos de disponibilização da base.

\textbf{Variáveis de deficiência.} As perguntas sobre deficiência,
Transtorno do Espectro Autista ou altas habilidades do professor e do
diretor ficam fora de todas as especificações. São dados de saúde, sem
relação com a pergunta de pesquisa.

\textbf{Cor/raça: duas especificações e uma auditoria.} A cor/raça declarada de
professores e diretores é dado pessoal sensível nos termos do
art.~5\textordmasculine{}, inciso~II, da Lei Geral de Proteção de Dados
(LGPD, Lei n\textordmasculine{} 13.709/2018). Essa qualificação não impede
o uso em pesquisa com microdados públicos e anonimizados, porque a própria
lei prevê o tratamento para a realização de estudos por órgão de pesquisa
(art.~7\textordmasculine{}, IV, e art.~11, II, ``c''). Pesquisas sobre
desigualdade educacional usam variáveis de raça legitimamente, e esse uso
pode ser condição para tornar desigualdades visíveis. O trabalho trata a
cor/raça, por isso, em duas especificações estimadas com o mesmo protocolo. Na
especificação A, a principal, ela fica fora do conjunto de preditores. Na especificação
B, entram quatro variáveis: as proporções de turmas cujo professor se
declara preto ou pardo, branco ou de outra cor ou sem declaração, e a
cor/raça do diretor. Os resultados das duas especificações são apresentados lado
a lado no Apêndice~\ref{ap:cor-raca}.

A especificação A é a principal por dois motivos. O primeiro é de escopo: a
pergunta de pesquisa (Seção~\ref{sec:problema-pesquisa}) trata de
características \emph{modificáveis por política educacional} (formação,
prática pedagógica, gestão, clima escolar), e atributos adscritos de
professores e diretores não pertencem a essa classe. O segundo é o risco de
leitura indevida. Com a cor/raça do docente entre os preditores, a
ordenação de importâncias poderia ser lida, ou reutilizada em um
\textit{dashboard}, como diferença de qualidade atribuível a raça, quando
refletiria a distribuição desigual de docentes entre escolas de contextos
distintos, mecanismos que os modelos aqui estimados não separam
(Seção~\ref{sec:registros-inferenciais}). A exclusão não custa precisão. Na
validação cruzada do treino, o MAE do modelo linear é de 0,1705 sem
cor/raça e de 0,1703 com ela na trilha de proporção, e de 17,86 e 17,89
pontos na de proficiência.

Excluir a variável, porém, não remove a informação que ela carrega. Um
modelo LightGBM treinado só com os preditores da especificação A recupera,
em validação cruzada, a cor/raça do diretor (preto ou pardo, contra branco)
com área sob a curva ROC de 0,72 e a proporção de turmas com professor
preto ou pardo com R\textsuperscript{2} de 0,14. A cor/raça permanece no
trabalho, assim, como variável de \emph{auditoria}. Os erros fora da
amostra das especificações A e B são comparados entre grupos de escolas
definidos pelo INSE, pela cor/raça do diretor e pela composição racial dos
professores. A mudança da diferença de erro entre os grupos, de A para B,
é reportada com intervalo de confiança \textit{bootstrap} de 95\%. Esses
intervalos são indicativos, porque não incorporam a incerteza do
retreinamento dos modelos. A ausência da variável no modelo principal não
autoriza concluir que não há desigualdade racial no desempenho das escolas.

\section{Estratégia empírica: comparação de modelos em duas trilhas}
\label{sec:estrategia-empirica}

A estratégia empírica deste trabalho consiste em uma comparação entre
algoritmos de aprendizado supervisionado, estimados em duas trilhas de
regressão complementares no nível da escola (a da proporção de alunos
alfabetizados e a da proficiência média em Língua Portuguesa), ambas
sobre exatamente o mesmo conjunto de preditores, cada uma excluindo apenas
a variável-alvo não utilizada pela outra trilha. Em cada trilha concorrem
um \emph{modelo linear regularizado de referência} e três
\textit{ensembles} de árvore, todos contrastados com um \textit{baseline}
trivial que estabelece o piso de desempenho abaixo do qual nenhum modelo
pode ser considerado útil. O modelo linear de referência é regularizado por
Elastic Net nas duas trilhas: uma regressão linear na trilha de proficiência
e, na trilha de proporção, o \textit{fractional logit} de
\citeonline{papke1996fractional}, especificação canônica para desfechos
fracionários. Ele representa, na terminologia de
\citeonline{breiman2001two}, a ``cultura'' da modelagem estatística clássica
diante da ``cultura'' algorítmica dos \textit{ensembles} de árvore.
A Tabela~\ref{tab:trilhas} resume a estratégia.

\begin{table}[H]
    \centering
    \footnotesize
    \caption{Trilhas de modelagem, alvos, modelos e métricas de avaliação}
    \begin{tabular}{P{2.2cm}P{2.2cm}P{3.0cm}P{1.6cm}P{3.2cm}}
        \toprule
        \textbf{Trilha} & \textbf{Alvo (nível escola)} & \textbf{Modelos} & \textbf{Métrica principal} & \textbf{Métricas complementares} \\
        \midrule
        Proporção de alfabetizados & Proporção de alunos alfabetizados
        $\in [0,1]$ & \textit{Baseline}: média/mediana. Referência linear:
        \textit{fractional logit} com Elastic Net. \textit{Ensembles}: Random
        Forest, XGBoost, LightGBM & MAE & RMSE, R\textsuperscript{2},
        análise de resíduos, aderência das predições ao intervalo $[0,1]$ \\ \addlinespace
        Proficiência média & Proficiência média em LP (contínua) &
        \textit{Baseline}: média/mediana. Referência linear: Regressão
        Linear com Elastic Net. \textit{Ensembles}: Random Forest, XGBoost,
        LightGBM & MAE & RMSE, R\textsuperscript{2}, sMAPE, análise de
        resíduos \\
        \bottomrule
    \end{tabular}
    \vspace{0.1cm}
    \begin{flushleft}
    {\footnotesize Fonte: elaborada pelo autor.\\
    Nota: MAE (\textit{Mean Absolute Error}): erro médio absoluto. RMSE
    (\textit{Root Mean Squared Error}): raiz do erro quadrático médio.
    R\textsuperscript{2}: coeficiente de determinação. sMAPE (\textit{symmetric
    Mean Absolute Percentage Error}): erro percentual absoluto médio simétrico.}
    \end{flushleft}
    \label{tab:trilhas}
\end{table}

A métrica principal das duas trilhas é definida \emph{previamente} à
estimação, de modo a evitar que a escolha do critério de comparação seja
influenciada pelos resultados observados. Adota-se o MAE (\emph{Mean
Absolute Error}, ou Erro Médio Absoluto, a média das distâncias absolutas
entre valor previsto e valor observado) em ambas: na
trilha de proporção, ele é lido diretamente em pontos percentuais de
alfabetização, a unidade em que as metas públicas são formuladas. Na trilha
de proficiência média, é expresso na própria escala do Saeb e é menos
sensível a observações extremas que o RMSE. As demais métricas são
reportadas como complementares, para caracterizar o comportamento dos
modelos, não para arbitrar a comparação. Na trilha de proporção, o MAE é
reportado em duas versões (não ponderada, cada escola com o mesmo peso,
coerente com a unidade de decisão, e ponderada pelo número de alunos
avaliados $n_e$), uma vez que escolas pequenas produzem proporções
mecanicamente mais voláteis e as duas leituras respondem a perguntas de
gestão distintas.

\subsection{Trilha da proporção de alfabetizados}

Na primeira trilha, o alvo é a proporção $Y_e^{\text{alf}} \in [0,1]$
definida na Equação~\ref{eq:prop-alfabetizado}. O modelo de referência
fracionário especifica a esperança condicional por meio de uma função de
ligação logística \cite{papke1996fractional}:

\begin{equation}
    \mathrm{E}\left[Y_e^{\text{alf}} \mid X_e\right] = \frac{1}{1 + e^{-(\beta_0 + \boldsymbol{\beta}^\top X_e)}},
    \label{eq:fraclogit}
\end{equation}

em que $X_e$ é o vetor de preditores da escola $e$. Essa especificação
garante, por construção, predições no intervalo unitário e acomoda sem
transformação ad hoc as escolas com proporções exatamente iguais a 0 ou 1
(presentes na base), situação em que a alternativa paramétrica da
regressão beta \cite{ferrari2004beta} tem verossimilhança indefinida, razão
pela qual esta é citada como alternativa, mas não adotada. A estimação
segue a quase-verossimilhança binomial com ligação logística proposta por
\citeonline{papke1996fractional}, com o número de alunos avaliados $n_e$
como peso de precisão e com a mesma penalidade Elastic Net da
Equação~\ref{eq:elasticnet}, necessária com centenas de preditores e 483
escolas de treino. Como o objetivo é preditivo, não se estimam erros-padrão
dos coeficientes. Os três \textit{ensembles} de árvore
(Random Forest \cite{breiman2001random}, XGBoost e LightGBM,
implementações do \textit{gradient boosting},
\citeonline{friedman2001greedy}) produzem, de forma geral, uma predição
aditiva sobre um conjunto de $M$ árvores de regressão $f_m$:

\begin{equation}
    \hat{y}_e = \sum_{m=1}^{M} f_m(X_e), \qquad f_m \in \mathcal{F},
    \label{eq:ensemble}
\end{equation}

em que $\mathcal{F}$ é o espaço de árvores de regressão. A diferença
essencial entre os três algoritmos está no procedimento de ajuste de cada
$f_m$: amostragem aleatória independente com agregação por média
(\textit{bagging}) no Random Forest, versus ajuste sequencial que corrige o
erro residual das árvores anteriores (\textit{boosting}) no XGBoost e no
LightGBM. Um detalhe de implementação relevante para o alvo fracionário
está documentado no código: as funções-objetivo empregadas devem aceitar
rótulos contínuos em $[0,1]$ (\texttt{reg:logistic} no XGBoost e
\texttt{cross\_entropy} no LightGBM). Os objetivos de classificação
binária ficam excluídos, porque exigem rótulos em $\{0,1\}$ e falhariam
silenciosamente com proporções.

\subsection{Trilha da proficiência média}

Na segunda trilha, o alvo passa a ser a proficiência média da escola, e o
modelo de referência é a Regressão Linear regularizada por Elastic Net, que
minimiza

\begin{equation}
    \min_{\boldsymbol{\beta}} \; \frac{1}{2n} \sum_{e=1}^{n} \left(y_e - \boldsymbol{\beta}^\top X_e\right)^2 + \lambda \left[\alpha \lVert \boldsymbol{\beta} \rVert_1 + \frac{1-\alpha}{2} \lVert \boldsymbol{\beta} \rVert_2^2\right],
    \label{eq:elasticnet}
\end{equation}

em que $\lambda$ controla a intensidade geral da penalização e $\alpha \in
[0,1]$ pondera entre a penalidade L1 (Lasso) e L2 (Ridge). Os mesmos três
\textit{ensembles} de árvore da primeira trilha são reestimados para este
alvo, seguindo a mesma forma geral da Equação~\ref{eq:ensemble}.

\subsection{Protocolo experimental}
\label{sec:protocolo}

Como a partição é feita por escola, o vazamento entre colegas de turma fica
excluído (Seção~\ref{sec:implicacoes-agregacao}). O protocolo abaixo, comum
às duas trilhas e a todos os modelos, separa estritamente decisão de
avaliação:

\begin{enumerate}
    \item \textbf{Separação treino/teste final por escola}: a partição é
    feita sobre as linhas-escola, na proporção aproximada de 80/20, com
    semente aleatória fixa documentada no código e estratificação pela
    distribuição do desfecho e pela região. A mesma partição vale nas duas
    trilhas, o que permite comparar diretamente a importância das
    variáveis entre elas;
    \item \textbf{Busca de hiperparâmetros restrita ao treino}: toda a busca
    é realizada exclusivamente sobre o conjunto de treino. O conjunto de
    teste final não participa de nenhuma decisão de modelagem: nem de
    seleção de hiperparâmetro, nem de escolha de modelo;
    \item \textbf{Validação cruzada}: a busca emprega validação cruzada de
    5 dobras sobre as escolas do conjunto de treino. A dependência
    residual entre escolas de um mesmo município
    (Seção~\ref{sec:implicacoes-agregacao}) é checada com uma validação
    cruzada agrupada por município (\texttt{GroupKFold});
    \item \textbf{Pré-processamento dentro de cada dobra}: a imputação
    (mediana para variáveis numéricas e ordinais, categoria ``não
    informado'' para nominais), a codificação \textit{one-hot} das
    variáveis nominais e a padronização usada pelos modelos lineares são
    ajustadas \emph{dentro} de cada dobra, a partir apenas dos dados de
    treinamento daquela dobra, e então aplicadas à dobra de validação. O
    mesmo pré-processamento vale para todos os modelos. A única etapa
    aplicada antes da partição é a exclusão de variáveis com mais de 30\%
    de ausência, calculada sobre toda a população, que usa apenas a
    ausência dos preditores e nenhuma informação dos desfechos;
    \item \textbf{Avaliação única no teste final}: concluída a seleção, o
    modelo escolhido de cada trilha é avaliado \emph{uma única vez} sobre o
    conjunto de teste. O valor assim obtido é a estimativa de desempenho
    reportada. Reavaliações sucessivas no teste após ajustes seriam,
    na prática, uma forma indireta de ajustar o modelo a ele;
    \item \textbf{Explicabilidade sobre o modelo selecionado}: o SHAP é
    aplicado ao modelo selecionado de cada trilha e às escolas do
    conjunto de teste, conforme a Seção~\ref{sec:xai}.
\end{enumerate}

Além do desempenho absoluto, cada modelo é comparado ao \textit{baseline}
trivial de sua trilha (média ou mediana do desfecho no treino). Assim, o
ganho informacional atribuível à modelagem fica explícito, além do valor da
métrica. Na trilha de proporção, verifica-se ainda a aderência
das predições ao intervalo $[0,1]$. Ela é garantida por construção no
\textit{fractional logit}, no Random Forest, cujas predições são médias de
proporções observadas, e no XGBoost e no LightGBM, estimados nessa trilha
com perda logística. Predições fora do intervalo, se ocorressem, seriam
truncadas e reportadas.

\section{Seleção do modelo final e testes de significância}
\label{sec:testes-significancia}

Para evitar que diferenças pequenas e possivelmente aleatórias entre
métricas sejam interpretadas como superioridade real de um algoritmo, a
escolha do modelo explicado em cada trilha segue uma regra definida antes
da estimação e aplicada exclusivamente ao conjunto de treino. O melhor
\textit{ensemble} (menor MAE de validação cruzada) é confrontado com o
modelo linear de referência nas \emph{mesmas} dobras de uma validação
cruzada de 5 dobras repetida cinco vezes. Para cada dobra $k$, calcula-se
a diferença $d_k = \text{MAE}_{\text{linear},k} - \text{MAE}_{\text{ensemble},k}$,
e o intervalo de confiança de 95\% da média de $d_k$ é obtido pelo teste $t$
corrigido de \citeonline{nadeau2003inference}, que ajusta a variância
estimada pela dependência entre dobras que compartilham observações de
treino:
\begin{equation}
    \widehat{\mathrm{Var}}(\bar{d}) = \left(\frac{1}{J} + \frac{n_{\text{val}}}{n_{\text{trein}}}\right) s^2_d,
    \label{eq:nadeau-bengio}
\end{equation}
em que $J$ é o número de dobras avaliadas e $s^2_d$ a variância amostral das
diferenças. O \textit{ensemble} é selecionado somente se (i) o limite
inferior do intervalo for positivo e (ii) o ganho relativo médio,
$\bar{d}/\overline{\text{MAE}}_{\text{linear}}$, for de pelo menos 2\%,
margem que representa a menor melhoria considerada relevante na prática.
Caso contrário, seleciona-se o modelo linear, mais simples e
intrinsecamente interpretável, em coerência com a objeção de
\citeonline{rudin2019stop} discutida na Seção~\ref{sec:xai-revisao}.

O conjunto de teste não participa dessa decisão. Concluída a seleção, ele é
utilizado uma única vez para reportar o desempenho do modelo escolhido, com
intervalo de confiança \textit{bootstrap} de 95\% para o MAE, e, de forma
descritiva, a diferença de erro entre o melhor \textit{ensemble} e o modelo
linear nas mesmas escolas, acompanhada do teste não paramétrico de
\textbf{Wilcoxon para amostras pareadas} sobre os erros absolutos. Com
121 escolas de teste, esse teste tem poder limitado para detectar
diferenças pequenas. Por isso, a ausência de significância não é tomada
como evidência de equivalência entre os modelos, e a decisão segue a regra
de seleção definida no treino.

Na trilha de proficiência média, reporta-se ainda, como métrica
complementar, o erro percentual absoluto médio simétrico,
$\text{sMAPE} = \frac{100}{n}\sum_{e} \frac{2\,|y_e - \hat{y}_e|}{|y_e| + |\hat{y}_e|}$,
que expressa o erro em termos relativos à escala do desfecho. Na trilha de
proporção, o sMAPE não é utilizado porque se torna instável quando a
proporção observada é igual ou próxima de zero.

\textbf{Perda de quantil.} A previsão também é avaliada fora da média, nos
quantis $\tau = 0{,}10$, $0{,}25$, $0{,}50$, $0{,}75$ e $0{,}90$ da
distribuição condicional de cada desfecho. Para cada $\tau$, calcula-se no
conjunto de teste a perda média de quantil (\textit{mean pinball loss})
\begin{equation}
    L_\tau = \frac{1}{n}\sum_{e} \rho_\tau\left(y_e - \hat{q}_{\tau,e}\right),
    \qquad \rho_\tau(u) = u\,\bigl(\tau - \mathrm{I}(u < 0)\bigr),
    \label{eq:pinball}
\end{equation}
em que $\hat{q}_{\tau,e}$ é o quantil previsto para a escola $e$ e
$\mathrm{I}(\cdot)$ é a função indicadora. A função $\rho_\tau$ pesa por
$\tau$ os erros acima do quantil e por $1-\tau$ os erros abaixo dele, e é a
mesma que a regressão quantílica minimiza \cite{koenker2005quantile}.
Concorrem, em cada $\tau$, a regressão quantílica linear com penalidade L1 e
as versões do XGBoost e do LightGBM com perda de quantil, contra um
\textit{baseline} igual ao quantil empírico do treino. O Random Forest fica
fora dessa comparação porque a implementação utilizada não estima quantis.
Cada modelo é calibrado só no treino, separadamente para cada $\tau$, com a
mesma validação cruzada da comparação principal e com a própria perda de
quantil como critério. Na trilha de proporção, todos os modelos usam o peso
de precisão $n_e$. As diferenças em relação ao \textit{baseline} são
reportadas com intervalo de confiança \textit{bootstrap} pareado de 95\%
sobre as escolas de teste. Como cada quantil é ajustado à parte, reporta-se
também a fração de escolas cujos quantis previstos saem fora de ordem. A
análise é complementar e não entra na regra de seleção do modelo final.

\section{Explicabilidade dos modelos (XAI)}
\label{sec:xai}

A explicabilidade concentra-se no \textbf{modelo selecionado de cada
trilha}, sem interpretar em detalhe todos os modelos estimados.
A razão é substantiva: importâncias produzidas por modelos que não são
adotados não informam sobre o fenômeno nem sobre o instrumento final, e
multiplicar explicações de modelos concorrentes convida à seleção
\textit{a posteriori} da narrativa mais conveniente.

O método prioritário é o SHAP, que segue a abordagem unificada de
\citeonline{lundberg2017shap}: para cada escola $e$, a predição do
modelo é decomposta em uma soma de contribuições aditivas e individualmente
atribuíveis a cada variável,

\begin{equation}
    f(X_e) = \phi_0 + \sum_{j=1}^{p} \phi_{e,j},
    \label{eq:shap}
\end{equation}

em que $\phi_0$ é o valor esperado da predição sobre o conjunto de treino e
$\phi_{e,j}$ é o valor de SHAP da $j$-ésima variável para a escola $e$.
Os valores são calculados sobre as \textbf{escolas do conjunto de
teste}, de modo que as explicações se refiram ao comportamento do modelo em
dados que não participaram de sua estimação. O explicador
\texttt{TreeExplainer} é utilizado quando o modelo selecionado for um
\textit{ensemble} de árvore, por computar valores de SHAP exatos em tempo
polinomial explorando a estrutura interna das árvores, conforme o algoritmo
TreeSHAP de \citeonline{lundberg2020tree}. O \texttt{LinearExplainer} é
utilizado caso o modelo selecionado seja o linear de referência.

A partir dos valores de SHAP do modelo selecionado, este trabalho produz
três níveis de análise. O primeiro é a importância agregada por variável,
resumida em gráficos de resumo (\textit{summary plots}). O segundo é a
importância agregada por grupo temático, obtida somando o valor absoluto
médio de SHAP de todas as variáveis do mesmo tema
(Seção~\ref{sec:variaveis}). Ela é mais informativa para a discussão de
política pública do que uma lista de centenas de variáveis isoladas, por
permitir perguntas do tipo ``o grupo `clima escolar' pesa mais que o grupo
`formação do professor'?''. O terceiro são gráficos de dependência
(\textit{dependence plots}) para as variáveis de maior importância
agregada, que mostram não linearidades na relação entre o valor da
variável e seu impacto na predição.

\textbf{Reagrupamento de categorias.} A codificação \textit{one-hot} de uma
variável nominal a fragmenta em várias colunas binárias, cada uma recebendo
seu próprio valor de SHAP. Reportar essas colunas isoladamente subestima
sistematicamente a importância da variável original e a torna incomparável
com variáveis numéricas, que permanecem em uma única coluna. Por isso, na
apresentação das importâncias, as colunas geradas por \textit{one-hot} de
uma mesma variável são \textbf{reagrupadas}, somando-se as contribuições
absolutas de suas categorias antes de qualquer ordenação ou agregação
temática.

\textbf{Estabilidade do ranking.} A ordenação dos temas por importância é
ela própria uma estimativa sujeita a variação amostral. Para não
interpretar como achado uma flutuação de estimação, a partição entre treino
e teste é refeita com 20 sementes aleatórias distintas. Em cada repetição,
o modelo selecionado é reajustado com os mesmos hiperparâmetros, e o SHAP é
recalculado no novo conjunto de teste. Reportam-se, para cada tema, a
posição média, o desvio-padrão da posição e a porcentagem de repetições em
que ele figura entre os três primeiros. Apenas ordenações estáveis entre
repetições são objeto de interpretação substantiva.

\textbf{LIME como checagem local.} O LIME \cite{ribeiro2016lime} é
aplicado a seis escolas do conjunto de teste, escolhidas por regra definida
antes da análise. São elas a escola de menor resíduo absoluto (caso
típico), a de maior resíduo positivo (acima do esperado), a de maior
resíduo negativo (abaixo do esperado), a de desfecho mais próximo da
mediana (caso limítrofe) e as de menor e de maior desfecho. Para cada
escola, o LIME é executado com dez sementes aleatórias, e suas
contribuições são reagrupadas na variável original, como no SHAP. Três
medidas resumem a comparação. A fidelidade local do LIME é o
R\textsuperscript{2} do modelo substituto na vizinhança da escola. A
estabilidade do LIME é o índice de Jaccard médio entre os conjuntos das dez
variáveis mais importantes de cada par de execuções. A concordância com o
SHAP da mesma escola é medida pelo índice de Jaccard entre as dez variáveis
mais importantes de cada método, pela correlação de Spearman entre as
contribuições e pela concordância de sinal nas variáveis em comum. As
divergências entre os dois métodos são tratadas como resultado a discutir.
A análise é complementar, e os objetivos do trabalho se apoiam no SHAP.

A interpretação de SHAP e LIME segue a regra de redação da
Seção~\ref{sec:registros-inferenciais}.

\section{Checagem de robustez}
\label{sec:robustez}

A restrição do escopo principal à rede pública, imposta pela
disponibilidade do INSE (Seção~\ref{sec:vies-inse}), poderia alterar as
conclusões sobre a importância relativa dos grupos temáticos. A
dificuldade de avaliar essa possibilidade é que a restrição altera
\emph{dois} fatores ao mesmo tempo: a \textbf{população analisada} (escolas
da rede pública apenas \textit{versus} universo público e privado) e o
\textbf{conjunto de preditores disponível} (com \textit{versus} sem o INSE e
as demais variáveis de \texttt{TS\_ESCOLA}). Uma comparação direta entre o
escopo principal e uma execução no universo sem INSE confundiria os dois
efeitos: qualquer diferença observada poderia ser atribuída tanto à mudança
de população quanto à retirada de um preditor reconhecidamente forte, sem
que se possa isolar a origem.

Por isso, a checagem de robustez adota um desenho que varia \textbf{um fator
de cada vez}, comparando três execuções de cada trilha sob o mesmo
protocolo da Seção~\ref{sec:protocolo}:

\begin{enumerate}
    \item \textbf{Execução de referência}. População: escolas da rede
    pública. Preditores: conjunto completo, incluindo o INSE. É o escopo
    principal do trabalho;
    \item \textbf{Variação do conjunto de preditores}. População: rede
    pública (a mesma da referência). Preditores: excluindo o INSE e as
    demais variáveis derivadas de \texttt{TS\_ESCOLA}. Comparada à
    referência, aproxima a contribuição do INSE, mantida constante a
    população, sabendo que preditores correlacionados podem compensar parte
    da sua retirada;
    \item \textbf{Variação da população}. População: universo de escolas
    públicas e privadas. Preditores: o mesmo conjunto reduzido da execução
    anterior. Comparada à execução 2, aproxima o efeito da ampliação da
    população, mantido constante o conjunto de preditores.
\end{enumerate}

A essas três execuções soma-se uma \textbf{execução de sensibilidade à
unidade de agregação}, em que as duas trilhas são repetidas com a
\emph{turma} como unidade, sob o mesmo protocolo (razão na
Seção~\ref{sec:agregacao}). A convergência das conclusões
entre as duas unidades indica que elas não dependem fortemente da escolha
da unidade de agregação, e uma divergência é reportada e discutida.

A leitura conjunta das três execuções principais permite aproximar
separadamente as diferenças observadas à retirada do INSE
(1~\textit{vs.}~2) e à mudança de população (2~\textit{vs.}~3). A análise privilegia a estabilidade do
\emph{ordenamento} dos grupos temáticos entre execuções, mais informativa
que a variação do valor absoluto da métrica, que responde também a
alterações na distribuição do desfecho e no tamanho amostral entre as
populações comparadas.

\textbf{Valor incremental dos blocos de variáveis.} Para responder
diretamente ao que os questionários contextuais acrescentam além do nível
socioeconômico e da geografia (Hipóteses 1 a 4), o modelo selecionado em cada
trilha é reestimado com os mesmos hiperparâmetros e nas mesmas dobras de uma
validação cruzada de 5 dobras repetida cinco vezes, nas especificações a
seguir.
\begin{enumerate}
    \item Somente o INSE.
    \item INSE e geografia.
    \item INSE, geografia e o bloco do professor.
    \item INSE, geografia e o bloco do diretor.
    \item Conjunto completo, usado como referência.
    \item Conjunto completo sem as duas variáveis de \texttt{TS\_ESCOLA} (o INSE
    e o percentual de docentes com formação adequada).
    \item Somente os blocos de professor e de diretor.
    \item Conjunto completo sem um bloco por vez (professor, diretor ou
    geografia), em três especificações.
    \item INSE, geografia e um índice por tema de cada questionário, no
    lugar dos itens. O índice é o primeiro componente principal das
    variáveis do tema, padronizadas e ajustadas só no treino de cada dobra,
    o que resulta em 19 índices do professor e 25 do diretor. Essa
    especificação foi declarada antes da execução final e estimada depois
    dela.
\end{enumerate}
O bloco de cada questionário inclui o indicador de escola sem aquele
questionário. Cada especificação é comparada ao conjunto completo e à
especificação com INSE e geografia pela diferença de MAE nas mesmas dobras,
com intervalo de confiança de 95\%
obtido pelo teste $t$ corrigido de \citeonline{nadeau2003inference}. Como os
hiperparâmetros são os escolhidos para o conjunto completo, a comparação
tende a favorecê-lo, o que torna conservadora qualquer conclusão de que um
bloco não acrescenta poder preditivo. As especificações que retiram um bloco
por vez do conjunto completo (8) e a que retira o INSE (6) aplicam a regra de
decisão das hipóteses (Seção~\ref{sec:hipoteses}). As demais descrevem o
ganho acumulado de cada bloco a partir do INSE e da geografia.

